\documentclass[12pt,a4paper]{article}
\usepackage[T1]{fontenc}
\usepackage[utf8]{inputenc}
\usepackage[ngerman]{babel}
\usepackage{amsmath,amssymb,amsthm,mathtools}
\usepackage{geometry}
\usepackage{fancyhdr}
\usepackage{titlesec}
\usepackage{enumitem}
% fraktur package removed - using mathfrak from amssymb instead
\usepackage{xfrac}

\geometry{margin=2.5cm}

\pagestyle{fancy}
\fancyhf{}
\fancyhead[C]{\nouppercase{\leftmark}}
\fancyhead[R]{\thepage}
\renewcommand{\headrulewidth}{0.4pt}

\newtheoremstyle{weber}% name
  {\topsep}% Space above
  {\topsep}% Space below
  {\normalfont}% Body font
  {}% Indent amount
  {\bfseries}% Theorem head font
  {.}% Punctuation after theorem head
  {.5em}% Space after theorem head
  {}% Theorem head spec (can be left empty, meaning ‘normal’)

\theoremstyle{weber}
\newtheorem*{satz}{Satz}
\newtheorem*{theorem}{Theorem}

\DeclareMathOperator{\Spur}{Spur}
\DeclareMathOperator{\Norm}{Norm}

\title{Lehrbuch der Algebra --- Band~1, Teil~13}
\author{Heinrich Weber}
\date{}

\begin{document}
\maketitle
\thispagestyle{empty}

\setcounter{section}{157}

\section*{\S.~158. Grundideal.}
\addcontentsline{toc}{section}{\S.~158. Grundideal}
\markboth{\S.~158. Grundideal}{\S.~158. Grundideal}

Dies widerspricht aber dem Satze 5.\ des vorigen Paragraphen,
nach dem $\tau$ nach einem Primzahlmodul $p$ keiner Congruenz von
niedrigerem als $n^{\mathrm{ten}}$ Grade gen"ugen kann.

Demnach ergiebt sich aus~(3), dass die Grundzahl $\varDelta$ des
K"orpers, vom Vorzeichen abgesehen, die absolute Norm
der Function $F'(\tau)$ ist, dass also
\begin{equation}\label{eq:5}
\pm\varDelta = N_a[F'(\tau)]
\end{equation}
zu setzen ist.

Wenn nun statt der $\omega_i$ eine andere Basis $\omega_i'$ von $\mathfrak{o}$ zu Grunde
gelegt wird, so tritt an Stelle von $\tau$ eine andere Form
\begin{equation}\label{eq:6}
\tau' = \omega_1't_1 + \omega_2't_2 + \dots + \omega_n't_n,
\end{equation}
wof"ur wir auch setzen k"onnen
\begin{equation}\label{eq:7}
\tau' = \omega_1t_1' + \omega_2t_2' + \dots + \omega_nt_n',
\end{equation}
und darin sind die $\omega_i'$ mit den $\omega_i$ durch eine lineare Substitution
mit der Determinante $\pm 1$ verbunden~(\S.~145):
\begin{equation}\label{eq:8}
(\omega_1', \omega_2', \dots, \omega_n') = C\,(\omega_1, \omega_2, \dots, \omega_n).
\end{equation}

F"uhrt man diese Substitution in~(6) aus, und ordnet nach
$\omega_1, \omega_2, \dots, \omega_n$, so ergiebt sich
\begin{equation}\label{eq:9}
(t_1', t_2', \dots, t_n') = C_1\,(t_1, t_2, \dots, t_n),
\end{equation}
wenn $C_1$ die transponirte Substitution von $C$ ist~(\S.~37). Bildet
man nun die Function $F(t)$ f"ur die Function $\tau'$, so mag sich $F_1(t)$
ergeben; die Ableitung f"ur $t = \tau'$ sei $F_1'(\tau')$. Setzen wir nun
\[
F'(\tau) = \varPsi(t_1, t_2, \dots, t_n),
\]
so ist $\varPsi$ eine ganze Function in $R$, und es ergiebt sich wegen~(7)
\[
F_1'(\tau') = \varPsi(t_1', t_2', \dots, t_n').
\]

Da die Substitution~(9) umkehrbar ist, so folgt hieraus, dass
die beiden Functionale $F'(\tau)$ und $F_1'(\tau')$ gegenseitig durch ein-
ander theilbar sind~(\S.~143), und dass sie mithin associirt sind.

Die Function $F'(\tau)$, deren absolute Norm gleich dem ab-
soluten Werthe der Grundzahl ist, nennen wir daher das \emph{Grund-\allowbreak functional}, und das durch $F'(\tau)$ repr"asentirte Ideal das \emph{Grundideal} des K"orpers~$\varOmega$. Nun sei $\mathfrak{p}^e$ die h"ochste Potenz des
Primideals $\mathfrak{p}$, die in $p$ aufgeht, und $e \leqq 1$, dann k"onnen wir
nach 5.\ des vorigen Paragraphen
\begin{equation}\label{eq:10}
F(t) \equiv P(t)^e\,\varPhi(t) \pmod{\mathfrak{p}}
\end{equation}
setzen, worin $\varPhi(t)$ eine ganze Function in $R$ von der Beschaffen-
heit ist, dass $\varPhi(\tau)$ nicht durch $\mathfrak{p}$ theilbar ist. Aus~(10) aber
folgt, indem wir die Ableitung nach $t$ bilden, was offenbar ge-
statlet ist:
\[
F'(t) \equiv e\,[P(t)]^{e-1}\,P'(t)\,\varPhi(t) + [P(t)]^e\,\varPhi'(t) \pmod{\mathfrak{p}}.
\]

Setzen wir hierin $t = \tau$, so geht $P'(t)$ in eine durch $\mathfrak{p}$
untheilbare Form $P'(\tau)$ "uber (was wir schon im Beweis von
\S.~157, 4.\ gezeigt und benutzt haben) und wir erhalten
\begin{equation}\label{eq:11}
F'(\tau) = e\,P(\tau)^{e-1}\,P'(\tau)\,F_1(\tau) \pmod{\mathfrak{p}^e}.
\end{equation}

Nun ist (nach \S.~157, 4.) $P(\tau)$ durch $\mathfrak{p}$, aber nicht durch $\mathfrak{p}^2$
theilbar, $P'(\tau)$ und $\varPhi(\tau)$ sind durch $\mathfrak{p}$ nicht theilbar, und so
giebt uns also die Formel~(11) den Beweis des folgenden Satzes:

\begin{satz}
Ist $\mathfrak{p}$ ein beliebiges Primideal, $p$ die durch $\mathfrak{p}$ theil-
bare nat"urliche Primzahl, und $\mathfrak{p}^e$ die h"ochste in
$p$ aufgehende Potenz von $\mathfrak{p}$, so ist die Grundform
$F'(\tau)$ allemal theilbar durch $\mathfrak{p}^{e-1}$; ist ferner der
Exponent $e$ nicht theilbar durch $\mathfrak{p}$, so ist $F'(\tau)$ nicht
theilbar durch $\mathfrak{p}^e$; ist aber $e$ theilbar durch $\mathfrak{p}$,
so ist $F'(\tau)$ theilbar durch $\mathfrak{p}^e$ und vielleicht durch
noch h"ohere Potenzen von~$\mathfrak{p}$.
\end{satz}

Wenn $e$ gr"osser als $1$ ist, so ist hiernach $F'(\tau)$ durch $\mathfrak{p}$
theilbar, und folglich ist $N_a[F'(\tau)]$ und also auch die Grundzahl~$\varDelta$
durch $p$ theilbar.

Wenn aber alle Primideale $\mathfrak{p}$ nur in erster Potenz in $p$ auf-
gehen, so ist $F'(\tau)$ relativ prim zu~$p$. Zerlegt man also $F'(\tau)$
in seine Primfactoren, so kommt darunter keiner vor, dessen
Norm eine Potenz von $p$ ist, folglich ist auch $N_a[F'(\tau)]$ und $\varDelta$
durch $p$ nicht theilbar. Daraus folgt dann der Satz:

\begin{satz}
Eine nat"urliche Primzahl $p$ ist dann und nur
dann im K"orper $\varOmega$ durch das Quadrat eines Prim-
factors theilbar, wenn $p$ in der Grundzahl von~$\varOmega$
aufgeht.
\end{satz}

Aus diesen Betrachtungen k"onnen wir noch andere wichtige
Schl"usse ziehen.

Wenn wir das System der linearen Gleichungen~(1) in Bezug
auf $\omega_1, \omega_2, \dots, \omega_n$ aufl"osen, so erhalten wir Ausdr"ucke mit dem
Nenner $U$, der, wie wir gesehen haben, eine Einheit ist. Sub-
stituirt man diese Ausdr"ucke in
\begin{equation}\label{eq:12}
\omega = x_1\omega_1 + x_2\omega_2 + \dots + x_n\omega_n,
\end{equation}
worin die $x_1, x_2, \dots, x_n$ ganze rationale Zahlen oder ganze
rationale Functionale sind, so erh"alt man einen Ausdruck von
der Form
\begin{equation}\label{eq:13}
\omega = A_0 + A_1\tau + A_2\tau^2 + \dots + A_{n-1}\tau^{n-1},
\end{equation}
worin die $A_0, A_1, \dots, A_{n-1}$ gleichfalls ganze rationale Func-
tionale bedeuten. Nach \S.~145, 2.\ wird aber in dieser Form jede
ganze Zahl und jedes ganze Functional des K"orpers $\varOmega$ dar-
gestellt, und wir k"onnen daher den Satz aussprechen:

\begin{satz}
Die Potenzen
\[
1,\; \tau,\; \tau^2,\; \dots,\; \tau^{n-1}
\]
bilden eine Basis der ganzen Functionale des
K"orpers~$\varOmega$.
\end{satz}

Statt der Potenzen von $\tau$ k"onnen auch die Functionen
\[
F_0, F_1, F_2, \dots, F_{n-1}
\]
als Elemente der Basis eingef"uhrt werden, die durch die Glei-
chung
\begin{equation}\label{eq:14}
\frac{F(t)}{t-\tau} = t^{n-1}F_0(\tau) + t^{n-2}F_1(\tau) + \dots + F_{n-1}(\tau)
\end{equation}
definirt sind, die, wie schon im \S.~4 des ersten Bandes gezeigt
ist, wenn
\[
F(t) = t^n + a_1t^{n-1} + \dots + a_n
\]
ist, den Ausdruck haben:
\begin{equation}\label{eq:15}
\begin{aligned}
F_0(t) &= 1, \\
F_1(t) &= t + a_1, \\
F_2(t) &= t^2 + a_1t + a_2, \\
&\;\vdotswithin{=} \\
F_{n-1}(t) &= t^{n-1} + a_1t^{n-2} + a_2t^{n-3} + \dots + a_{n-1},
\end{aligned}
\end{equation}
so dass die Potenzen $1, t, t^2, \dots, t^{n-1}$ ohne Nenner durch
$F_0, F_1, \dots, F_{n-1}$ dargestellt werden k"onnen, und dass jede ganze
Zahl oder jedes ganze Functional $\omega$ auch den Ausdruck erh"alt:
\begin{equation}\label{eq:16}
\omega = B_0F_0(\tau) + B_1F_1(\tau) + \dots + B_{n-1}F_{n-1}(\tau),
\end{equation}
dessen Coefficienten $B_0, B_1, \dots, B_{n-1}$ ganze rationale Func-
tionale sind.

Betrachten wir jetzt die ganze Function von~$t$:
\[
S\,\frac{F(t)}{(t-\tau)\,F'(\tau)},
\]
worin $S$ das Zeichen f"ur die in \S.~134 erkl"arte Spur ist, so
ergiebt sich, dass sie den Werth~$1$ erh"alt f"ur $t = \tau_1, \tau_2, \dots, \tau_n$
und dass sie also, da sie nur vom $(n-1)^{\mathrm{ten}}$ Grade ist, identisch
$= 1$ sein muss. Daraus folgt nach~(14):
\[
1 = t^{n-1}S\,\frac{F_0(\tau)}{F'(\tau)} + t^{n-2}S\,\frac{F_1(\tau)}{F'(\tau)} + \dots + S\,\frac{F_{n-1}(\tau)}{F'(\tau)}
\]
oder
\begin{equation}\label{eq:17}
\begin{gathered}
S\,\frac{F_\nu(\tau)}{F'(\tau)} = 0, \quad S\,\frac{F_{n-1}(\tau)}{F'(\tau)} = 1, \\
\nu = 0, 1, \dots, n-2.
\end{gathered}
\end{equation}

Hieraus ergiebt sich mit Benutzung von~(16):
\begin{equation}\label{eq:18}
S\,\frac{\omega}{F'(\tau)} = B_{n-1},
\end{equation}
oder der Satz:

\begin{satz}
Ist $\omega$ eine ganze Zahl oder ein ganzes Func-
tional des K"orpers $\varOmega$, so ist die Spur von $\dfrac{\omega}{F'(\tau)}$
ein ganzes rationales~Functional.
\end{satz}

\bigskip\bigskip
\centerline{$*\;*\;*$}
\bigskip\bigskip

\section*{Neunzehnter Abschnitt.}
\addcontentsline{toc}{section}{Neunzehnter Abschnitt. Beziehungen eines K"orpers auf seine Theiler}

\begin{center}
\Large\bfseries
Beziehungen eines K"orpers auf seine Theiler.
\end{center}
\bigskip

\section*{\S.~159. Relativnormen.}
\addcontentsline{toc}{section}{\S.~159. Relativnormen}
\markboth{\S.~159. Relativnormen}{\S.~159. Relativnormen}

Wir wollen in diesem Abschnitte die Modificationen be-
trachten, die in den Definitionen und S"atzen der Theorie der
algebraischen Zahlen eintreten, wenn an Stelle des absoluten
Rationalit"atsbereiches, d.~h.\ des K"orpers der rationalen
Zahlen, ein beliebiger algebraischer Zahlk"orper gesetzt wird. Es
werden sich dabei einige wichtige Erg"anzungen auch f"ur die
allgemeine Theorie der Primfactoren ergeben, die in einer Reihe
von Dedekind und Hilbert aufgestellter S"atze ihren Aus-
druck finden\footnote{Dedekind, "Ueber die Discriminanten endlicher K"orper, Abhand-
lungen der G"ottinger Gesellschaft der Wissenschaften (1882). ,,Zur Theorie
der Ideale``, Nachrichten der Gesellschaft der Wissenschaften zu G"ottingen,
1894, Nr.~4. Hilbert, ,,Grundz"uge einer Theorie des Galois'schen Zahl-
k"orpers``. Ebend. 1894, Nr.~3.}.

Es sei also $R$ ein algebraischer Zahlk"orper $m^{\mathrm{ten}}$ Grades und
\begin{equation}\label{eq:s159-1}
f(\varTheta) = \varTheta^n + \alpha_1\varTheta^{n-1} + \dots + \alpha_n = 0
\end{equation}
eine in $R$ rationale und irreducible Gleichung $n^{\mathrm{ten}}$ Grades. Der
K"orper $\varOmega = R(\varTheta)$ ist ein algebraischer K"orper "uber $R$, der in
Bezug auf $R$ vom $n^{\mathrm{ten}}$ Grade ist (Bd.~I, \S.~142). Im absoluten
Rationalit"atsbereiche ist $\varOmega$ vom Grade $mn$, und wenn $\varrho$ eine
Primitivzahl des K"orpers $R$ ist, so ist
\begin{equation}\label{eq:s159-2}
\begin{array}{ll}
\varTheta^s\varrho^t, & s = 0, 1, \dots, n-1 \\
                 & t = 0, 1, \dots, m-1
\end{array}
\end{equation}
eine Basis des K"orpers $\varOmega$; denn wegen der Irreducibilit"at der
Gleichung~(1) kann zwischen den $mn$ Zahlen~(2) keine lineare
Relation mit rationalen Zahlencoefficienten bestehen.

Jede Zahl $\omega$ des K"orpers $\varOmega$ kann als ganze Function
$(n-1)^{\mathrm{ten}}$ Grades von $\varTheta$ mit Coefficienten in $R$ dargestellt werden,
und jede solche Zahl $\omega$ gen"ugt einer in $R$ irreduciblen Gleichung
h"ochstens vom $n^{\mathrm{ten}}$ Grade; $\omega$ ist dann und nur dann eine ganze
Zahl, wenn diese Gleichung, nachdem der Coefficient der h"och-
sten Potenz auf $1$ gebracht ist, ganze Zahlen in $R$ zu Coeffi-
cienten hat.

Den Inbegriff aller ganzen Zahlen in $\varOmega$ bezeichnen wir, wie
fr"uher, mit $\mathfrak{o}$.

Wenn wir, ohne die Zahlen des K"orpers $R$ zu ver"andern,
f"ur $\varTheta$ die s"ammtlichen Wurzeln $\varTheta_1, \varTheta_2, \dots, \varTheta_n$ der Gleichung~(1)
nehmen, so erhalten wir $n$ K"orper
\begin{equation}\label{eq:s159-3}
\varOmega_1 = R(\varTheta_1),\; \varOmega_2 = R(\varTheta_2),\; \dots,\; \varOmega_n = R(\varTheta_n),
\end{equation}
die in Bezug auf den K"orper $R$ conjugirt sind.

Sind $\omega_1, \omega_2, \dots, \omega_n$ entsprechende Zahlen dieser K"orper, so
heisst das Product
\begin{equation}\label{eq:s159-4}
\mathfrak{N}_R(\omega) = \omega_1\omega_2\dots\omega_n
\end{equation}
die in Bezug auf $R$ genommene Partialnorm (\emph{Relativ-\allowbreak norm}) der Zahl~$\omega$.

Eine Zahl $\eta$ des K"orpers $R$ hat eine in diesem K"orper
genommene gew"ohnliche Norm $N_R(\eta)$, die eine rationale Zahl
ist. Die Relativnorm $\mathfrak{N}_R(\omega)$ ist nun eine solche Zahl $\eta$, und
wenn wir ihre Norm im K"orper $R$ nehmen, so erhalten wir die
Totalnorm der Zahl $\omega$ im K"orper der rationalen Zahlen:
\begin{equation}\label{eq:s159-5}
N_\varOmega(\omega) = N_R\,\mathfrak{N}_R(\omega).
\end{equation}

Dies ergiebt sich, wenn man $\omega$ durch die Potenzen~(2) von
$\varTheta$ und $\varrho$ ausdr"uckt, sodann f"ur $\varrho$ die $m$ conjugirten Werthe,
und zu jedem dieser $\varrho$ die nach $R$ conjugirten Werthe $\varTheta$ setzt.

In demselben Sinne haben nun auch die Functionale des
K"orpers $\varOmega$ und die durch diese repr"asentirten Ideale ihre
Partialnormen. Die Partialnorm eines Functionals in $\varOmega$ ist
ein Functional in $R$, und demnach ist die Partialnorm eines
Ideals in $\varOmega$ ein Ideal in~$R$.

Von den conjugirten Idealen k"onnen auch zwei oder mehrere
einander gleich sein, und zwar kann dies auch dann eintreten,
wenn die die Ideale repr"asentirenden conjugirten Functionale
nicht identisch sind, wenn sie nur associirt sind. Wenn man
daher nur das Product aller von einander verschiedenen unter
den conjugirten Idealen nimmt, so braucht dies keineswegs ein
Ideal in $R$ zu sein.

Was hier von den Normen ausgef"uhrt ist, gilt auch von
den anderen symmetrischen Functionen, insbesondere von den
Spuren. Wir nennen die Summe
\begin{equation}\label{eq:s159-6}
\mathfrak{S}_R(\omega) = \omega_1 + \omega_2 + \dots + \omega_n
\end{equation}
die Partialspur (oder \emph{Relativspur}) von $\omega$ in Bezug
auf~$R$. Sie ist eine Zahl oder ein Functional in $R$, und wir er-
halten f"ur die Totalspur
\begin{equation}\label{eq:s159-7}
S_\varOmega(\omega) = S_R\,\mathfrak{S}_R(\omega),
\end{equation}
worin die Bedeutung der Zeichen unmittelbar verst"andlich ist.

Eine besonders wichtige Beziehung ergiebt sich aber durch
die Betrachtung der Grundideale.

Es sei $\tau$ eine Basisform von $\varOmega$ (also eine Linearform von
$mn$ Variablen, \S.~146) und
\[
F(t) = N_\varOmega\,(t-\tau)
\]
die irreducible Function $mn^{\mathrm{ten}}$ Grades (im absoluten Rationalit"ats-
bereiche), deren Wurzel $\tau$ ist. Dann haben wir im \S.~158 das
durch das Functional $F'(\tau)$ definirte Ideal als das \emph{Grundideal}
des K"orpers $\varOmega$ bezeichnet.

Ist nun $\sigma$ eine Basisform von $R$ und
\[
\psi(t) = N_R\,(t-\sigma),
\]
so ist durch $\psi'(\sigma)$ das Grundideal des K"orpers $R$ definirt.

Aus der Function $F(t)$ spaltet sich aber im K"orper $R$ ein
irreducibler Factor $n^{\mathrm{ten}}$ Grades ab,
\[
f(t) = \mathfrak{N}_R\,(t-\tau),
\]
der f"ur $t = \tau$ verschwindet. Wir definiren durch $f'(\tau)$ das
Partial-Grundideal von $\varOmega$ in Bezug auf $R$, und beweisen
nun den Satz:
\begin{equation}\label{eq:s159-8}
F'(\tau) = \varepsilon\,f'(\tau)\,\psi'(\sigma),
\end{equation}
worin $\varepsilon$ eine Einheit ist.

Um ihn zu beweisen, setzen wir nach der im \S.~158 an-
gewandten Bezeichnung
\begin{align}
\frac{F(t)}{t-\tau} &= t^{mn-1}F_0(\tau) + t^{mn-2}F_1(\tau) + \dots + F_{mn-1}(\tau), \label{eq:s159-9}\\
\frac{f(t)}{t-\tau} &= t^{n-1}f_0(\tau) + t^{n-2}f_1(\tau) + \dots + f_{n-1}(\tau), \label{eq:s159-10}\\
\frac{\psi(t)}{t-\sigma} &= t^{m-1}\psi_0(\sigma) + t^{m-2}\psi_1(\sigma) + \dots + \psi_{m-1}(\sigma). \label{eq:s159-11}
\end{align}

Hierin sind $F_\nu, \psi_\nu$ ganze Functionale des absoluten Ratio-
nalit"atsbereiches, w"ahrend $f_\nu$ ganze Functionale in $R$ sind. Die
Ausdr"ucke f"ur $F_\nu, f_\nu, \psi_\nu$ ergeben sich aus den Formeln \S.~158,~(15).
Auf diese Functionale kann man dieselben Schl"usse anwenden,
die zu dem Theorem \S.~158, 4.\ gef"uhrt haben, und es ergiebt sich:
\begin{equation}\label{eq:s159-12}
\begin{gathered}
\mathfrak{S}_R\,\frac{f_\nu(\tau)}{f'(\tau)} = 0, \\
\mathfrak{S}_R\,\frac{f_{n-1}(\tau)}{f'(\tau)} = 1,
\end{gathered}
\qquad \nu = 0, 1, \dots, n-2.
\end{equation}

Wenn man nun die Functionen $F_\nu(t)$, wenn sie von h"oherem
als $(n-1)^{\mathrm{ten}}$ Grade sind, durch $f(t)$ dividirt und den Rest der
Division nimmt, dann in diesem Reste die Potenzen von~$t$ [nach
\S.~158,~(15), auf $f_\nu$ angewandt], durch die Functionen $f_\nu(t)$ aus-
dr"uckt, so ergiebt sich die Darstellung
\begin{equation}\label{eq:s159-13}
F_\nu(\tau) = \alpha_0f_0 + \alpha_1f_1 + \dots + \alpha_{n-1}f_{n-1},
\end{equation}
worin die $\alpha_0, \alpha_1, \dots, \alpha_{n-1}$ ganze Functionale in $R$ sind. Diese
Darstellung gilt f"ur jede Wurzel $\tau$ von $f(t) = 0$, d.~h.\ f"ur
$\tau = \tau_1, \tau_2, \dots, \tau_n$. Hieraus ergiebt sich aber nach~(12):
\[
\mathfrak{S}_R\left(\frac{F_\nu(\tau)}{f'(\tau)}\right) = \alpha_{n-1}.
\]

Dividirt man hier durch $\psi'(\sigma)$ und bildet die Spur $S_R$, so
erh"alt man aus~(7) mit Anwendung des Satzes~4., \S.~158:
\begin{equation}\label{eq:s159-14}
S_\varOmega\left(\frac{F_\nu(\tau)}{f'(\tau)\,\psi'(\sigma)}\right) = a_\nu,
\end{equation}
worin $a_\nu$ ein ganzes rationales Functional bedeutet.

Nach~(9) ist nun f"ur jede Wurzel $\tau_i$ von $F(t) = 0$ mit
Ausnahme von $\tau_i = \tau$
\[
\tau^{nm-1}F_0(\tau_i) + \tau^{nm-2}F_1(\tau_i) + \dots + F_{nm-1}(\tau_i) = 0,
\]
und
\[
\tau^{nm-1}F_0(\tau) + \tau^{nm-2}F_1(\tau) + \dots + F_{nm-1}(\tau) = F'(\tau),
\]
so dass sich, wenn man~(14) mit $\tau^{nm-\nu-1}$ multiplicirt und die
Summe in Bezug auf $\nu$ nimmt,
\begin{equation}\label{eq:s159-15}
\frac{F'(\tau)}{f'(\tau)\,\psi'(\sigma)} = a_0\tau^{nm-1} + a_1\tau^{nm-2} + \dots
\end{equation}
ergiebt. Es ist daher $F'(\tau)$ durch $f'(\tau)\,\psi'(\sigma)$ theilbar.

Andererseits ist $f_\nu(\tau)\,\psi_\mu(\sigma)$ ein ganzes Functional des K"or-
pers $\varOmega$, und folglich nach dem vorhin schon angewandten Satze
\S.~158,~4.:
\begin{equation}\label{eq:s159-16}
S_\varOmega\left(\frac{f_\nu(\tau)\,\psi_\mu(\sigma)}{F'(\tau)}\right) = b_{\mu,\nu}
\quad
\begin{array}{l}
\mu = 0, 1, \dots, m-1 \\
\nu = 0, 1, \dots, n-1
\end{array}
\end{equation}
ein ganzes rationales Functional. Multiplicirt man diese Glei-
chung mit $\tau^{n-\nu-1}\sigma^{m-\mu-1}$ und nimmt die Summe "uber alle $\mu$
und $\nu$, so folgt, wie vorhin, da f"ur zwei von $\sigma$, $\tau$ verschiedene
Wurzeln $\sigma_k, \tau_i$ von $\psi = 0, f = 0$
\begin{gather*}
\sigma^{m-1}\psi_0(\sigma_k) + \sigma^{m-2}\psi_1(\sigma_k) + \dots + \psi_{m-1}(\sigma_k) = 0,\\
\tau^{n-1}f_0(\tau_i) + \tau^{n-2}f_1(\tau_i) + \dots + f_{n-1}(\tau_i) = 0,
\intertext{und}
\sigma^{m-1}\psi_0(\sigma) + \sigma^{m-2}\psi_1(\sigma) + \dots + \psi_{m-1}(\sigma) = \psi'(\sigma),\\
\tau^{n-1}f_0(\tau) + \tau^{n-2}f_1(\tau) + \dots + f_{n-1}(\tau) = f'(\tau)
\end{gather*}
ist, die Gleichung
\begin{equation}\label{eq:s159-17}
\frac{f'(\tau)\,\psi'(\sigma)}{F'(\tau)} = \sum^{\mu,\nu} b_{\mu,\nu}\,\tau^{n-\nu-1}\sigma^{m-\mu-1}.
\end{equation}

Also ist auch $f'(\tau)\,\psi'(\sigma)$ durch $F'(\tau)$ theilbar und unser
Theorem~(8) daher bewiesen.

Bezeichnen wir mit $G_\varOmega, G_R$ die Grundideale der K"orper $\varOmega, R$,
mit $\mathfrak{G}_R$ das relative Grundideal von $\varOmega$ in Bezug auf $R$, so kann
man dem bewiesenen Theorem auch den Ausdruck geben:
\begin{equation}\label{eq:s159-18}
G_\varOmega = \mathfrak{G}_R\,G_R.
\end{equation}

Es ist klar, wie sich dieses Theorem verallgemeinern l"asst:
Es sei
\[
\varOmega_1, \varOmega_2, \varOmega_3, \dots
\]
eine Reihe von K"orpern, deren jeder die folgenden als Theiler
enth"alt. Es sei ferner $\mathfrak{G}_{i,k}$ das Grundideal von $\varOmega_i$ in Bezug auf
$\varOmega_k$ ($k > i$), dann ist
\begin{equation}\label{eq:s159-19}
\mathfrak{G}_{i,k} = \mathfrak{G}_{i,i+1}\;\mathfrak{G}_{i+1,i+2}\dots\mathfrak{G}_{k-1,k}.
\end{equation}

Die tiefere Bedeutung dieser Grundideale wird sich in einem
der n"achsten Paragraphen noch deutlicher herausstellen.

Wir wollen aber schon hier auf eine merkw"urdige Er-
scheinung aufmerksam machen, dass es n"amlich vorkommen kann,
dass $f'(\tau)$ eine Einheit ist, was nach dem Minkowski'schen
Satze nicht m"oglich ist, wenn $R$ der K"orper der rationalen
Zahlen ist. In diesen F"allen, von denen die Theorie der com-
plexen Multiplication der elliptischen Functionen Beispiele giebt,
ist das Partial-Grundideal $\mathfrak{G}_R$ gleich $1$ zu setzen.

Wenn die $n$ K"orper~(3) mit einander identisch sind, so
heisst $\varOmega$ ein \emph{Normalk"orper} in Bezug auf $R$ (relativ normal,
wenn $R$ nicht der absolute Rationalit"atsbereich ist). Die Gesamt-
heit der $n$ Substitutionen $(\varTheta, \varTheta_\nu)$, die den Uebergang der Zahlen
des K"orpers $\varOmega$ zu den conjugirten Zahlen vermitteln, bilden
eine Gruppe vom Grade $n$, wie wir schon im \S.~147 des ersten
Bandes gesehen haben, die nichts anderes ist, als die \emph{Galois'sche
Gruppe} der Gleichung~(1) im Rationalit"atsbereiche $R$, die wir
hier die \emph{Gruppe des K"orpers $\varOmega$} (in Bezug auf $R$) nennen.

Wenn insbesondere diese Gruppe commutativ ist, so ist $\varOmega$
ein \emph{Abel'scher K"orper} in Bezug auf $R$ (ein relativ Abel'scher
K"orper). In diesem Falle ist~(1) eine Abel'sche Gleichung im
Rationalit"atsbereiche~$R$.

\newpage
\section*{\S.~160. Primideale im relativ normalen K"orper.}
\addcontentsline{toc}{section}{\S.~160. Primideale im relativ normalen K"orper}
\markboth{\S.~160. Primideale im relativ normalen K"orper}{\S.~160. Primideale im relativ normalen K"orper}

Wir nehmen jetzt an, dass der K"orper $\varOmega$ normal sei in
Bezug auf den beliebigen K"orper $R$, und untersuchen unter
dieser Voraussetzung die Primideale des K"orpers~$\varOmega$.

Die Gruppe von $\varOmega$ in Bezug auf $R$ sei $\varPhi$, und die Sub-
stitutionen von $\varPhi$ m"ogen mit $\varphi, \varphi_1, \dots$ bezeichnet sein. Wir
bezeichnen ferner mit
\begin{equation}\label{eq:s160-1}
\omega \mid \varphi
\end{equation}
die Zahl, die durch die Substitution $\varphi$ aus $\omega$ hervor-
geht, so dass $\omega$ und $\omega \mid \varphi$ in $\varOmega$ enthalten sind.

Ebenso wie die Zahlen $\omega$ gehen auch alle Functionale und
damit zugleich alle Ideale $\mathfrak{a}$ des K"orpers $\varOmega$ durch eine Sub-
stitution $\varphi$ in bestimmte Functionale und Ideale $\mathfrak{a} \mid \varphi$ "uber, und
wenn $\omega$ eine durch $\mathfrak{a}$ theilbare ganze Zahl ist, so ist $\omega \mid \varphi$ durch
$\mathfrak{a} \mid \varphi$ theilbar.

Wenn $\mathfrak{a}$ irgend ein Ideal in $\varOmega$ ist, so giebt es gewisse Sub-
stitutionen $\psi$ in $\varPhi$, die der Bedingung
\begin{equation}\label{eq:s160-2}
\mathfrak{a} \mid \psi = \mathfrak{a}
\end{equation}
gen"ugen, darunter immer die identische Substitution, und diese
Substitutionen $\psi$ bilden eine Gruppe $\varPsi$, die ein Theiler von $\varPhi$
ist. Wir nennen $\varPsi$ die zum Ideal $\mathfrak{a}$ geh"orige Gruppe, und
wir sagen auch, $\mathfrak{a}$ geh"ort zu der Gruppe~$\varPsi$.

Da man in jeder Gleichung zwischen Zahlen und Func-
tionalen in $\varOmega$ alle Substitutionen der Gruppe $\varPhi$ ausf"uhren darf,
wobei die Gr"ossen des K"orpers $R$ unge"andert bleiben, ohne dass
die Gleichung zu bestehen aufh"ort, so folgt, dass Einheiten,
ganze Functionale, associirte Functionale, durch einander theil-
bare Functionale diese Eigenschaften nicht verlieren, wenn irgend
eine der Substitutionen von $\varPhi$ ausgef"uhrt wird. Wir heben den
Satz hervor:

\begin{satz}
Ist $\mathfrak{p}$ ein Primideal, so sind auch alle mit $\mathfrak{p}$ in
Bezug auf $R$ conjugirten Ideale $\mathfrak{p} \mid \varphi$ Primideale.
Ist $\mathfrak{a}$ durch irgend eine Potenz von $\mathfrak{p}$ theilbar,
so ist $\mathfrak{a} \mid \varphi$ durch die gleiche Potenz von $\mathfrak{p} \mid \varphi$,
theilbar.
\end{satz}

Ist $\mathfrak{p}$ ein Primideal in $\varOmega$, so giebt es eine und nur eine
durch $\mathfrak{p}$ theilbare nat"urliche Primzahl~$p$.

Zerlegt man diese in ihre Primfactoren in $R$, so muss einer
dieser Primfactoren, den wir mit $\mathfrak{P}$ bezeichnen, durch $\mathfrak{p}$ theil-
bar sein.

Ist dann $\mathfrak{A}$ irgend ein Ideal in $R$ und zugleich durch $\mathfrak{p}$
theilbar, so ist der gr"osste gemeinschaftliche Theiler von $\mathfrak{P}$
und $\mathfrak{A}$, der nach \S.~139 auch in $R$ enthalten ist, durch $\mathfrak{p}$ theil-
bar, und ist also gewiss keine Einheit. $\mathfrak{P}$ und $\mathfrak{A}$ sind also nicht
relativ prim, und $\mathfrak{A}$ ist folglich durch das Primideal $\mathfrak{P}$ theilbar.
Ist $\mathfrak{A}$ auch ein Primideal, so m"ussen $\mathfrak{A}$ und $\mathfrak{P}$ identisch sein.
Damit ist, mit R"ucksicht auf~1., bewiesen:

\begin{satz}
Ist $\mathfrak{p}$ ein Primideal in $\varOmega$, so giebt es ein und nur
ein Primideal $\mathfrak{P}$ in $R$, das durch $\mathfrak{p}$ theilbar ist,
und jedes durch $\mathfrak{p}$ theilbare Ideal in $R$ ist zu-
gleich durch $\mathfrak{P}$ theilbar. Das Primideal $\mathfrak{P}$ ist
zugleich durch die s"ammtlichen zu $\mathfrak{p}$ conjugirten
Primideale $\mathfrak{p} \mid \varphi$ theilbar.
\end{satz}

Die Partialnorm $\mathfrak{N}_R(\mathfrak{p})$ ist durch $\mathfrak{p}$, und folglich als Ideal in
$R$ durch $\mathfrak{P}$ theilbar. Sie kann aber auch nach dem Satze~1.
durch kein von $\mathfrak{P}$ verschiedenes Ideal in $R$ theilbar sein. Denn
ist $\mathfrak{P}'$ irgend ein Primtheiler von $\mathfrak{N}_R(\mathfrak{p})$, so ist $\mathfrak{P}'$ wenigstens durch
einen der Primfactoren $\mathfrak{p} \mid \varphi$ theilbar und mithin gleich $\mathfrak{P}$. Folglich
ist die Partialnorm von $\mathfrak{p}$ eine Potenz von~$\mathfrak{P}$. Wir setzen
\begin{equation}\label{eq:s160-3}
\mathfrak{N}_R(\mathfrak{p}) = \mathfrak{P}^f,
\end{equation}
und nennen $f$ den \emph{Grad} von $\mathfrak{p}$ in Bezug auf den K"orper~$R$.

Die Norm von $\mathfrak{P}$ im K"orper $R$ ist gleichfalls eine Potenz
von $p$, die wir mit $P$ bezeichnen wollen, also
\begin{equation}\label{eq:s160-4}
N_R(\mathfrak{P}) = P.
\end{equation}

Es ergiebt sich dann f"ur die Totalnorm von $\mathfrak{p}$ im K"orper $\varOmega$
nach \S.~159,~(5)
\begin{equation}\label{eq:s160-5}
N_\varOmega(\mathfrak{p}) = P^f.
\end{equation}

Ist $\mathfrak{p}^g$ die h"ochste in $\mathfrak{P}$ aufgehende Potenz von $\mathfrak{p}$, so ist $\mathfrak{P}$
nach~1.\ auch durch die $g^{\mathrm{te}}$ Potenz aller mit $\mathfrak{p}$ conjugirten Prim-
factoren und durch keine h"ohere Potenz theilbar. Wenn nun
$\mathfrak{p}_1, \mathfrak{p}_2, \dots, \mathfrak{p}_e$ die von einander verschiedenen unter den mit
$\mathfrak{p}$ conjugirten Primidealen sind, so ist
\begin{equation}\label{eq:s160-6}
\mathfrak{P} = (\mathfrak{p}_1\,\mathfrak{p}_2\dots\mathfrak{p}_e)^g,
\end{equation}
und wenn man die Partialnorm auf beiden Seiten nimmt, und
wie fr"uher mit $n$ den Grad des K"orpers $\varOmega$ in Bezug auf $R$ be-
zeichnet, so dass die Partialnorm von $\mathfrak{P}$ gleich $\mathfrak{P}^n$ wird, so folgt
aus~(3):
\begin{equation}\label{eq:s160-7}
n = efg.
\end{equation}
Es sind also die nat"urlichen Zahlen $e, f, g$ Theiler von~$n$.

Wenn $\mathfrak{p}$ durch $\varphi_1$ in $\mathfrak{p}_1$ "ubergeht, wenn also
\[
\mathfrak{p}_1 = \mathfrak{p} \mid \varphi_1,
\]
so ist $\mathfrak{p} = \mathfrak{p}_1 \mid \varphi_1^{-1}$, und wenn $\varPsi$ die zu $\mathfrak{p}$ geh"orige Gruppe ist,
so ist auch f"ur jede in $\varPsi$ enthaltene Substitution $\psi$
\[
\mathfrak{p}_1 = \mathfrak{p} \mid \psi\varphi_1, \quad \mathfrak{p}_1 = \mathfrak{p}_1 \mid \varphi_1^{-1}\psi\varphi_1.
\]

Ist umgekehrt $\mathfrak{p}_1 \mid \varphi = \mathfrak{p}_1$, so folgt $\mathfrak{p} \mid \varphi_1\varphi\varphi_1^{-1} = \mathfrak{p}$, d.~h.\
$\varphi_1\varphi\varphi_1^{-1} = \psi$ oder $\varphi = \varphi_1^{-1}\psi\varphi_1$.

Es geh"ort daher $\mathfrak{p}_1$ zur Gruppe $\varphi_1^{-1}\varPsi\varphi_1$, und wenn
\[
\mathfrak{p}_1 = \mathfrak{p} \mid \varphi_1, \; \mathfrak{p}_2 = \mathfrak{p} \mid \varphi_2, \; \dots, \; \mathfrak{p}_e = \mathfrak{p} \mid \varphi_e
\]
ist, so ist
\begin{equation}\label{eq:s160-8}
\varPhi = \varPsi\varphi_1 + \varPsi\varphi_2 + \dots + \varPsi\varphi_e.
\end{equation}
$\varPsi$ ist also ein Theiler von $\varPhi$ vom Index $e$ und vom Grade~$gf$.

\newpage
\section*{\S.~161. Primitiveurzeln der Primideale.}
\addcontentsline{toc}{section}{\S.~161. Primitiveurzeln der Primideale}
\markboth{\S.~161. Primitiveurzeln der Primideale}{\S.~161. Primitiveurzeln der Primideale}

Aus den Formeln~(4), (5) des vorigen Paragraphen ergiebt
sich, mit R"ucksicht auf \S.~150, f"ur jede ganze Zahl $\eta$ des K"or-
pers~$R$
\[
\eta^P \equiv \eta \pmod{\mathfrak{P}},
\]
also auch
\begin{equation}\label{eq:s161-1}
\eta^P \equiv \eta \pmod{\mathfrak{p}},
\end{equation}
und f"ur jede Zahl $\omega$ in $\mathfrak{o}$
\begin{equation}\label{eq:s161-2}
\omega^{Pf} \equiv \omega \pmod{\mathfrak{p}}.
\end{equation}

Die Anzahl der nach $\mathfrak{P}$ incongruenten ganzen Zahlen in $R$
ist $P$, und die Anzahl der nach $\mathfrak{p}$ incongruenten Zahlen in $\varOmega$
ist~$P^f$.

Wir haben schon fr"uher~(\S.~150) gezeigt, dass es zu jedem
Primideal $\mathfrak{p}$ Primitivwurzeln $\gamma$ giebt, d.~h.\ Zahlen in $\varOmega$, die
der Congruenz
\begin{equation}\label{eq:s161-3}
\gamma^{Pf-1} \equiv 1 \pmod{\mathfrak{p}}
\end{equation}
gen"ugen, und zugleich die Eigenschaft haben, dass keine niedrigere
Potenz mit positiven Exponenten der Einheit congruent wird.
Dann ist jede durch $\mathfrak{p}$ nicht theilbare ganze Zahl in $\varOmega$ mit einer
und nur mit einer der Potenzen von $\gamma$
\begin{equation}\label{eq:s161-4}
1,\; \gamma,\; \gamma^2,\; \dots,\; \gamma^{Pf-2}
\end{equation}
nach dem Modul $\mathfrak{p}$ congruent.

Jede ganze Zahl $\varTheta$ in $\varOmega$ gen"ugt einer Gleichung h"ochstens
vom $n^{\mathrm{ten}}$ Grade, deren Coefficienten ganze Zahlen in $R$ sind.

Diese Gleichung ist zugleich eine Congruenz nach dem
\textbf{Modul} $\mathfrak{p}$, und unter allen solchen Congruenzen wird es eine von
m"oglichst niedrigem Grade
\[
F(\varTheta) \equiv 0 \pmod{\mathfrak{p}}
\]
geben, in der wir "uberdies den Coefficienten der h"ochsten Potenz
von $\varTheta$ gleich $1$ annehmen k"onnen. Denn ist der h"ochste (durch
$\mathfrak{p}$ und folglich durch $\mathfrak{P}$ untheilbare) Coefficient $\eta_0$, so k"onnen
wir immer der Congruenz
\[
\eta_0\eta \equiv 1 \pmod{\mathfrak{P}}
\]
durch ein ganzzahliges $\eta$ in $R$ gen"ugen, und haben dann nur
die Function $F$ mit diesem $\eta$ zu multipliciren. Wir k"onnen
also, wenn $t$ eine Variable ist, die Function $F$ in der Form an-
nehmen:
\[
F(t) = t^r + \alpha_1t^{r-1} + \alpha_2t^{r-2} + \dots + \alpha_r,
\]
deren Coefficienten $\alpha_1, \alpha_2, \dots, \alpha_r$ ganze Zahlen in $R$ sind. Durch
Division k"onnen wir dann f"ur jede andere ganze Function $F_1(t)$
den Quotienten $Q(t)$ und den Rest $\varphi(t)$ so bestimmen, dass
\begin{equation}\label{eq:s161-5}
F_1(t) = Q(t)\,F(t) + \varphi(t),
\end{equation}
worin
\begin{equation}\label{eq:s161-6}
\varphi(t) = \varrho_0 + \varrho_1t + \dots + \varrho_{r-1}t^{r-1}
\end{equation}
eine durch $F_1(t)$ eindeutig bestimmte Function $(r-1)^{\mathrm{ten}}$ Grades
ist, mit ganzzahligen Coefficienten~$\varrho$.

Setzen wir f"ur $\varTheta$ die Primitivwurzel $\gamma$ und f"ur $F_1(t)$ eine
Potenz von $t$, so ergiebt sich aus~(4), (5) und~(6), dass jede durch
$\mathfrak{p}$ nicht theilbare Zahl $\omega$ einer Congruenz
\begin{equation}\label{eq:s161-7}
\omega \equiv \varrho_0 + \varrho_1\gamma + \dots + \varrho_{r-1}\gamma^{r-1} \pmod{\mathfrak{p}}
\end{equation}
gen"ugt, und da die $\varrho$ auch $= 0$ sein k"onnen, so gilt dies auch
noch f"ur ein durch $\mathfrak{p}$ theilbares $\omega$.

Die Coefficienten $\varrho$ in dem Ausdrucke~(7) sind durch $\omega$
selbst nach dem Modul $\mathfrak{P}$ v"ollig bestimmt, da nach unserer
Voraussetzung $\gamma$ keiner Congruenz in $R$ von niedrigerem als dem
$r^{\mathrm{ten}}$ Grade gen"ugen soll, deren Coefficienten nicht alle durch $\mathfrak{P}$
theilbar sind. Folglich giebt es, da jeder der Coefficienten in~(7)
nur $P$ incongruente Werthe haben kann, $P^r$ und nicht mehr in-
congruenten Zahlen $\omega$, und daraus folgt, dass $\nu = f$ sein muss.

Aus einer Congruenz $F(\gamma) \equiv 0 \pmod{\mathfrak{p}}$ folgt nun aber
durch wiederholtes Potenziren mit R"ucksicht auf~(1):
\[
F(\gamma^P) \equiv 0,\; F(\gamma^{P^2}) \equiv 0,\; \dots \pmod{\mathfrak{p}},
\]
und wir haben also den Satz bewiesen~(\S.~150, 2.):

\begin{satz}
Eine Primitivwurzel $\gamma$ von $\mathfrak{p}$ gen"ugt nach
dem Modul $\mathfrak{p}$ einer Congruenz in $R$ vom $f^{\mathrm{ten}}$ Grade,
deren s"ammtliche Wurzeln
\[
\gamma,\; \gamma^P,\; \gamma^{P^2},\; \dots,\; \gamma^{P^{f-1}}
\]
sind.
\end{satz}

Diesem Satze kann man auch den Ausdruck geben:

Das Product
\begin{equation}\label{eq:s161-8}
(t-\gamma)\,(t-\gamma^P)\dots(t-\gamma^{P^{f-1}})
\end{equation}
ist nach dem Modul $\mathfrak{p}$ mit einer ganzen Func-
tion $F(t)$ in $R$ congruent.

\newpage
\section*{\S.~162. Das Partial-Grundideal.}
\addcontentsline{toc}{section}{\S.~162. Das Partial-Grundideal}
\markboth{\S.~162. Das Partial-Grundideal}{\S.~162. Das Partial-Grundideal}

Es m"oge jetzt
\begin{equation}\label{eq:s162-1}
\tau = t_1\omega_1 + t_2\omega_2 + \dots + t_{mn}\omega_{mn}
\end{equation}
eine Basisform von $\mathfrak{o}$ sein. Dann wird es gewisse Substitutionen
$\chi$ in $\varPhi$ geben, und darunter immer die identische, die der Con-
gruenz
\begin{equation}\label{eq:s162-2}
\tau \mid \chi \equiv \tau \pmod{\mathfrak{p}}
\end{equation}
gen"ugen. Alle diese Substitutionen bilden eine in $\varPhi$ enthaltene
Gruppe $X$, die auch ein Theiler der Gruppe $\varPsi$ ist, zu der $\mathfrak{p}$
geh"ort. Denn aus $\tau$ erh"alt man (nach \S.~146) eine Basisform
von $\mathfrak{p}$, wenn man f"ur die Variablen $t$ gewisse lineare Functionen
neuer Variablen mit ganzen rationalen Coefficienten setzt; wenn
also $\pi$ eine solche Basisform von $\mathfrak{p}$ ist, so folgt aus~(2), dass $\pi \mid \chi$
durch $\pi$ theilbar und daher auch $\mathfrak{p} \mid \chi = \mathfrak{p}$ ist.

Nun gen"ugt $\tau$ einer Gleichung $n^{\mathrm{ten}}$ Grades $f(t) = 0$, deren
Coefficenten ganze Functionen in $R$ mit den Variablen $t_1, t_2 \dots$
sind, und es ist, wenn $t$ eine neue Variable bedeutet, und
$\tau_1, \tau_2, \dots, \tau_n$ die in Bezug auf $R$ zu $\tau$ conjugirten Formen sind,
\begin{equation}\label{eq:s162-3}
f(t, \tau_1, \tau_2, \dots) = f(t) = (t-\tau_1)\,(t-\tau_2)\dots(t-\tau_n).
\end{equation}

Hieraus ergiebt sich nun, wenn wir wiederholt in die Po-
tenz $P$ erheben, und auf die Zahlencoefficienten von $f$ die Formel
\S.~161,~(1) anwenden:
\begin{equation}\label{eq:s162-4}
\begin{aligned}
f\,(t^P, t_1^P, t_2^P, \dots) &\equiv [f(t)]^P \\
&\equiv (t^P-\tau_1')\,(t^P-\tau_2')\dots(t^P-\tau_n') \pmod{\mathfrak{p}},
\end{aligned}
\end{equation}
wenn $\tau_1', \tau_2', \dots, \tau_n'$ dadurch aus $\tau_1, \tau_2, \dots, \tau_n$ abgeleitet sind, dass
die Variablen $t_1, t_2, \dots$ durch $t_1^P, t_2^P, \dots$ ersetzt sind.

Setzen wir nun $t = \tau$ in~(4), so verschwindet $f(t)$ und es
folgt, dass einer der Factoren des letzten Productes durch $\mathfrak{p}$
theilbar sein muss. Dies aber kann so ausgedr"uckt werden, dass
es in $\varPhi$ eine Substitution $\psi_0$ giebt, die der Bedingung
\begin{equation}\label{eq:s162-5}
\tau^P \equiv \tau' \mid \psi_0 \pmod{\mathfrak{p}}
\end{equation}
gen"ugt.

Aus $\tau$ entstehen alle ganzen Zahlen in $\varOmega$, wenn man f"ur
die Variablen ganze rationale Zahlen setzt. Wendet man dann
noch den Fermat'schen Lehrsatz f"ur rationale Zahlen an, so er-
kennt man, dass die Congruenzen~(2) und~(5) gleichbedeutend
sind mit den f"ur jede Zahl $\omega$ in $\mathfrak{o}$ g"ultigen Formeln
\begin{equation}\label{eq:s162-6}
\omega \mid \chi \equiv \omega, \quad \omega^P \equiv \omega \mid \psi_0 \pmod{\mathfrak{p}}.
\end{equation}

Aus der zweiten Congruenz~(6) folgt, dass, wenn $\omega$ durch $\mathfrak{p}$
theilbar ist, immer auch $\omega \mid \psi_0$ durch $\mathfrak{p}$ theilbar sein muss. Folg-
lich ist $\mathfrak{p}$ durch $\mathfrak{p} \mid \psi_0$ theilbar, und als Primideal $= \mathfrak{p} \mid \psi_0$. Daraus
folgt, dass $\psi_0$ in der Gruppe $\varPsi$ enthalten ist.

Wir verstehen jetzt unter $\gamma$ eine Primitivwurzel von $\mathfrak{p}$ und
wenden die am Schlusse des \S.~161 bewiesene Formel an:
\begin{equation}\label{eq:s162-7}
(t-\gamma)\,(t-\gamma^P)\dots(t-\gamma^{P^{f-1}}) \equiv F(t) \pmod{\mathfrak{p}},
\end{equation}
in der $F(t)$ eine ganze Function von $t$ in $R$ ist.

Daraus folgt
\[
F(\gamma) \equiv 0 \pmod{\mathfrak{p}},
\]
und wenn nun $\psi$ irgend eine Substitution aus $\varPsi$ ist:
\[
F(\gamma \mid \psi) \equiv 0 \pmod{\mathfrak{p}}.
\]

Daraus ergiebt sich aber, dass $\gamma \mid \psi$ mit einer der in~(7) vor-
kommenden Potenzen von $\gamma$ nach $\mathfrak{p}$ congruent sein muss, also
etwa
\begin{equation}\label{eq:s162-8}
\gamma^{P^r} \equiv \gamma \mid \psi \pmod{\mathfrak{p}}.
\end{equation}

Nun ist jede durch $\mathfrak{p}$ nicht theilbare Zahl $\omega$ in $\mathfrak{o}$ mit einer
Potenz von $\gamma$ congruent, und wenn man also~(8) zu dieser Potenz
erhebt, so folgt
\begin{equation}\label{eq:s162-9}
\omega^{P^r} \equiv \omega \mid \psi \pmod{\mathfrak{p}},
\end{equation}
und diese Congruenz gilt offenbar auch noch, wenn $\omega$ durch $\mathfrak{p}$
theilbar ist, also f"ur alle Zahlen in~$\mathfrak{o}$.

Wenn wir nun die zweite der Congruenzen~(6) $\nu$ mal nach
einander anwenden, so folgt
\begin{equation}\label{eq:s162-10}
\omega^{P^\nu} \equiv \omega \mid \psi_0^\nu \pmod{\mathfrak{p}},
\end{equation}
also
\begin{equation}\label{eq:s162-11}
\omega \mid \psi \equiv \omega \mid \psi_0^r \pmod{\mathfrak{p}},
\end{equation}
und wenn wir $\omega$ durch $\omega \mid \psi^{-1}$ ersetzen oder auf~(11) die Sub-
stitution $\psi^{-1}$ anwenden:
\begin{equation}\label{eq:s162-12}
\omega \equiv \omega \mid \psi^{-1}\psi_0^r \equiv \omega \mid \psi_0^r\psi^{-1} \pmod{\mathfrak{p}}.
\end{equation}
Es sind also sowohl $\psi^{-1}\psi_0^r$ als auch $\psi_0^r\psi^{-1}$ in $X$ enthalten,

voraus sich ergiebt, dass jede Substitution $\psi$ aus $\varPsi$ in einem
der Systeme (Nebengruppen)
\[
X,\; X\psi_0,\; X\psi_0^2,\; \dots
\]
enthalten ist.

Wenn $\mathfrak{p}$ vom Grade $f$ (in Bezug auf $R$) ist, so ist
\[
\omega^{Pf} \equiv \omega \pmod{\mathfrak{p}},
\]
und $Pf$ ist die niedrigste Potenz von $P$ mit positiven Expo-
nenten, die dieser Congruenz f"ur alle $\omega$ gen"ugt.

Setzen wir also $\nu = f$ in~(10), so folgt, dass $\psi_0^f$ in $X$ ent-
halten ist, dass aber keine Potenz von $\psi_0$ mit niedrigerem posi-
tiven Exponenten diese Eigenschaft hat. Die Nebengruppen
\[
X,\; X\psi_0,\; X\psi_0^2,\; \dots,\; X\psi_0^{f-1}
\]
sind also alle von einander verschieden, und es ergiebt sich:
\begin{equation}\label{eq:s162-13}
\varPsi = X + X\psi_0 + X\psi_0^2 + \dots + X\psi_0^{f-1}.
\end{equation}

Nach~(12) ist aber die Gesammtheit der Substitutionen $\psi_0^\nu X$
mit $X\psi_0^\nu$ identisch (f"ur jedes $\nu$), also
\begin{equation}\label{eq:s162-14}
\psi_0^\nu X = X\psi_0^\nu,
\end{equation}
hieraus folgt, dass $X$ ein Normaltheiler von $\varPsi$ ist.

Da, wie wir oben gesehen haben, $\varPsi$ vom Grade $gf$
ist, so ist $X$ vom Grade~$g$.

In der Gruppe $X$ giebt es nun solche Substitutionen $\chi_1$,
darunter immer die identische, f"ur die die Congruenz
\begin{equation}\label{eq:s162-15}
\tau \mid \chi_1 \equiv \tau \pmod{\mathfrak{p}^2}
\end{equation}
erf"ullt ist, und diese bilden eine Gruppe $X^{(1)}$, deren Grad mit $g_1$
bezeichnet sein mag.

Es ist zu beweisen, dass $g_1$ eine Potenz von $p$ ist.
Bezeichnen wir n"amlich mit $\pi$ ein durch $\mathfrak{p}$ theilbares Prim-
functional, mit $\alpha$ ein ganzes Functional, so folgt aus~(15):
\begin{equation}\label{eq:s162-16}
\tau \mid \chi_1 = \tau + \pi^2\alpha,
\end{equation}
und da ebenso nach~(15)
\[
\alpha \mid \chi_1 \equiv \alpha \pmod{\mathfrak{p}^2}, \quad (\pi \mid \chi_1)^2 \equiv \pi^2 \pmod{\mathfrak{p}^3}
\]
ist, so ergiebt sich aus~(16)
\[
\tau \mid \chi_1^2 \equiv \tau + 2\pi^2\alpha \pmod{\mathfrak{p}^3},
\]
und allgemein f"ur jeden positiven Exponenten $h$
\[
\tau \mid \chi_1^h \equiv \tau + h\pi^2\alpha \pmod{\mathfrak{p}^3}.
\]

Setzt man $h = p$, so ergiebt sich hieraus
\[
\tau \mid \chi_1^p \equiv \tau \pmod{\mathfrak{p}^3},
\]
und auf die gleiche Weise schliesst man hieraus f"ur jeden Ex-
ponenten $\nu$
\begin{equation}\label{eq:s162-17}
\tau \mid \chi_1^{p^\nu} \equiv \tau \pmod{\mathfrak{p}^{\nu+2}}.
\end{equation}
Um diese Formel allgemein zu beweisen, wendet man die
vollst"andige Induction an. Man nimmt~(17) als bewiesen an
und setzt
\[
\tau \mid \chi_1^{p^\nu} = \tau + \pi^{\nu+2}\alpha;
\]
daraus durch wiederholte Anwendung
\begin{align*}
\tau \mid \chi_1^{2p^\nu} &\equiv \tau + 2\pi^{\nu+2}\alpha \\
\tau \mid \chi_1^{3p^\nu} &\equiv \tau + 3\pi^{\nu+2}\alpha \pmod{\mathfrak{p}^{\nu+3}} \\
&\;\vdotswithin{\equiv} \\
\tau \mid \chi_1^{p^{\nu+1}} &\equiv \tau + p\pi^{\nu+2}\alpha,
\end{align*}
also die Formel~(17) f"ur $\nu + 1$.

Nun kann man $\nu$ immer so gross annehmen, dass von den
Differenzen
\[
\tau - \tau_1,\; \tau - \tau_2,\; \tau - \tau_3,\; \dots,
\]
wenn $\tau$ von $\tau_1, \tau_2, \tau_3, \dots$ verschieden ist, keine durch $\mathfrak{p}^{\nu-2}$ theil-
bar ist, und dann folgt aus~(17), dass $\chi_1^{p^\nu}$ die identische Sub-
stitution sein muss. Es ist hiernach der Grad eines jeden
Elementes $\chi_1$ eine Potenz von $p$, und folglich kann
nach dem Cauchy-Sylow'schen Satze~(\S.~29, I.) im Grade
von $X^{(1)}$ keine von $p$ verschiedene Primzahl aufgehen,
wie bewiesen werden sollte.

Wenn also $g$ nicht durch $p$ theilbar ist, dann ist sicher $X^{(1)}$
die Einheitsgruppe.

Man kann auf diese Weise fortfahren und eine Gruppe $X^{(2)}$
vom Grade $g_2$ bilden aus der Congruenz
\[
\tau \mid \chi_2 \equiv \tau \pmod{\mathfrak{p}^3},
\]
$X^{(2)}$ ist ein Theiler von $X^{(1)}$, also $g_2$ gleichfalls eine Potenz von $p$,
und so muss man schliesslich zur Einheitsgruppe gelangen.
Unter den auf einander folgenden Gruppen $X^{(1)}, X^{(2)}, \dots$ k"onnen
unter Umst"anden auch mehrere einander gleich sein\footnote{Hilbert nennt $\varPsi$ die Zerlegungsgruppe, $X$ die Tr"agheitsgruppe,
$X^{(1)}$ die Verzweigungsgruppe des Ideals $\mathfrak{p}$. In der oben erw"ahnten Abhand-
lung ist noch bewiesen, dass $g_1$ die h"ochste in $g$ aufgehende Potenz von $p$
ist, dass $X^{(1)}$ ein Normaltheiler von $X$ ist, und dass die ganze Gruppe $\varPsi$
metacyklisch ist.}.

\newpage
\section*{\S.~163. Die Ideale in den Theilern des K"orpers $\varOmega$.}
\addcontentsline{toc}{section}{\S.~163. Die Ideale in den Theilern des K"orpers $\varOmega$}
\markboth{\S.~163. Die Ideale in den Theilern des K"orpers $\varOmega$}{\S.~163. Die Ideale in den Theilern des K"orpers $\varOmega$}

Hiernach l"asst sich die h"ochste Potenz des Primideals $\mathfrak{p}$ an-
geben, die in dem Partialgrundideal $\mathfrak{G}_R$ aufgeht. Nach \S.~159
n"amlich ist dies Ideal definirt durch das Product
\[
f'(\tau) = (\tau-\tau_1)\,(\tau-\tau_2)\dots(\tau-\tau_{n-1}),
\]
worin $\tau_1, \tau_2, \dots, \tau_{n-1}$ die von $\tau$ verschiedenen unter den mit $\tau$
conjugirten Functionalen sind. Einer der Factoren dieses Pro-
ductes, $\tau - \tau \mid \varphi$, ist dann und nur dann durch $\mathfrak{p}$ theilbar, wenn
$\varphi$ zu $X$ geh"ort. Da nun die identische Substitution ausgeschlossen
ist, so folgt, dass $\mathfrak{G}$ den Factor $\mathfrak{p}$ mindestens $g-1$ mal enth"alt.

Es ist aber $\tau - \tau \mid \chi$ dann und nur dann durch $\mathfrak{p}^2$ theilbar,
wenn $\chi$ zu $X^{(1)}$ geh"ort, und folglich ist $\mathfrak{G}$ durch $\mathfrak{p}^{g-1+g_1-1}$
theilbar. Wir k"onnen so fortfahren und finden die genaue Potenz
von $\mathfrak{p}$, die in $\mathfrak{G}$ enthalten ist. Sie hat den Exponenten
\[
(g-1) + (g_1-1) + (g_2-1) + \dots
\]
Jede der Zahlen $g, g_1, g_2, \dots$ ist Theiler der vorangehenden,
und $g_1, g_2, \dots$ sind Potenzen von $p$, und die Reihe setzt sich so
lange fort, bis eine der Zahlen $g = 1$ geworden ist. Die Zu-
sammensetzung des Grundideals $\mathfrak{G}$ ist also vollkommen durch
die Zerlegung der Gruppe $X$ bestimmt.

Die Zahlen $g, g_1, g_2, \dots$ sind "uberdies f"ur alle conjugirten
Primideale $\mathfrak{p}_1, \mathfrak{p}_2, \dots, \mathfrak{p}_e$ dieselben, und es folgt also
\[
\mathfrak{G} = \varPi\,(\mathfrak{p}_1\,\mathfrak{p}_2\dots\mathfrak{p}_e)^{g-1+g_1-1+g_2-1\dots},
\]
wo sich das Productzeichen auf die Primfactoren aller Prim-
zahlen $p$ des K"orpers $\varOmega$ bezieht. Es brauchen dabei aber selbst-
verst"andlich nur die in endlicher Anzahl vorhandenen Prim-
zahlen $p$ ber"ucksichtigt zu werden, f"ur die $g > 1$ ist.

Die Partialnorm des Ideals $\mathfrak{G}$
\[
\mathfrak{D} = \mathfrak{N}_R(\mathfrak{G}) = \varPi\;\mathfrak{P}^{ef(g-1+g_1-1+g_2-1+\dots)}
\]
heisst die \emph{Partial-Discriminante} des K"orpers $\varOmega$ in Bezug auf den
K"orper $R$. Sie ist ein im K"orper $R$ gelegenes Ideal, die f"ur den
Fall, dass $R$ der K"orper der rationalen Zahlen ist, von einem
Einheitsfactor abgesehen, in die Grundzahl des K"orpers $\varOmega$ "uber-
geht~(\S.~158).

\newpage
Wenn die Gruppe $\varPhi$ des K"orpers $\varOmega$ einen Theiler $\varPhi'$ hat,
so geh"ort zu $\varPhi'$ ein K"orper $\varOmega'$, der ein Theiler von $\varOmega$ ist und
seinerseits $R$ als Theiler enth"alt. Die Zahlen von $\varOmega'$ bleiben
durch die Substitutionen $\varPhi'$ unge"andert, und $\varPhi'$ ist die Gruppe des
K"orpers $\varOmega$ in Bezug auf $\varOmega'$ (vergl.\ Bd.~I, \S.~155). Wir bezeichnen
den Grad von $\varPhi'$, der ein Theiler von $n$ ist, mit $n'$ und setzen
\begin{equation}\label{eq:s163-1}
n = m'n'.
\end{equation}

Es ist dann $n'$ der Grad von $\varOmega$ in Bezug auf $\varOmega'$, und $m'$
der Grad von $\varOmega'$ in Bezug auf~$R$.

Es ist dabei zu beachten, dass zwar $\varOmega$ ein Normalk"orper
auch in Bezug auf $\varOmega'$ ist, dass aber $\varOmega'$ im Allgemeinen kein
Normalk"orper in Bezug auf $R$ sein wird.

Die Primideale im K"orper $\varOmega'$ lassen sich nun vollst"andig aus
denen von $\varOmega$ ableiten.

Wenn $\mathfrak{p}$ irgend ein Primideal in $\varOmega$ ist, so giebt es ein und
nur ein durch $\mathfrak{p}$ theilbares Ideal $\mathfrak{p}'$ in $\varOmega'$, und $\mathfrak{p}'$ kann als
Theiler von $\mathfrak{P}$ durch keine anderen als die mit $\mathfrak{p}$ conjugirten Ideale
theilbar sein. Durchl"auft $\varphi'$ die Substitutionen von $\varPhi'$, so ist $\mathfrak{p}'$,
da es durch $\varphi'$ unge"andert bleibt, auch durch $\mathfrak{p} \mid \varphi'$ theilbar.

Wenn also $\mathfrak{p}_1'$ durch $\mathfrak{p} \mid \varphi_1$ theilbar ist, so ist es auch durch
$\mathfrak{p} \mid \psi\varphi_1\varphi'$ theilbar, wenn $\psi$ ein beliebiges Element der Gruppe $\varPsi$
(\S.~162) bedeutet.

Wenn daher $\varphi$ irgend ein Element des Systemes
\[
\varPhi_1 = \varPsi\varphi_1\varPhi'
\]
ist, so ist $\mathfrak{p}_1'$ durch $\mathfrak{p} \mid \varphi$ theilbar. Es kommt demnach jetzt die
Gruppenzerlegung in Betracht, die wir im \S.~28 dieses Bandes
auseinandergesetzt haben.

Wenn $\varphi_2$ ein nicht in $\varPhi_1$ enthaltenes Element aus $\varPhi$ ist, so
hat das System
\[
\varPhi_2 = \varPsi\varphi_2\varPhi'
\]
mit $\varPhi_1$ gar kein Element gemein. Giebt es noch ein Element $\varphi$
das weder in $\varPhi_1$ noch in $\varPhi_2$ vorkommt, so bildet man ebenso
\[
\varPhi_3 = \varPsi\varphi_3\varPhi',
\]
und f"ahrt so fort, bis die ganze Gruppe $\varPhi$ ersch"opft ist. Man
erh"alt so die \emph{Zerlegung}
\begin{equation}\label{eq:s163-2}
\varPhi = \varPhi_1 + \varPhi_2 + \varPhi_3 + \dots + \varPhi_{e'}.
\end{equation}

Den Grad eines dieser Systeme $\varPhi_r$ k"onnen wir nach \S.~28
bestimmen. Er ist gleich dem Producte des Grades von $\varPhi$
multiplicirt mit dem Index des Durchschnittes $\varPsi_r$ von $\varPhi'$ mit
$\varPsi_r = \varphi_r^{-1}\varPsi\varphi_r$, in Bezug auf $\varPhi'$, oder was dasselbe ist, gleich
dem Producte der Grade von $\varPsi$ und $\varPhi'$, getheilt durch den
Grad $h_r$ des Durchschnittes von $\varPhi'$ und $\varPsi_r$. Der Grad von $\varPhi_r$
ist also gleich $fgn' : h_r$.

Wenn wir jetzt mit $\varphi_r$ alle Elemente von $\varPhi_r$ bezeichnen, so
f"uhren alle Primideale $\mathfrak{p} \mid \varphi_r$ in $\varOmega$ zu demselben Primideale $\mathfrak{p}_r'$
in $\varOmega'$. Es ist aber noch zu zeigen, dass zwei verschiedene dieser
Systeme $\varPhi_1, \varPhi_2$ zu verschiedenen Primidealen $\mathfrak{p}_1', \mathfrak{p}_2'$ f"uhren.

Zun"achst ist ersichtlich, dass die s"ammtlichen $\mathfrak{p} \mid \varphi_2$ von den
$\mathfrak{p} \mid \varphi_1$ verschieden sind. Denn wenn $\mathfrak{p} \mid \varphi_1 = \mathfrak{p} \mid \varphi_2$ ist, so ist
auch $\mathfrak{p} = \mathfrak{p} \mid \varphi_2\varphi_1^{-1}$, also $\varphi_2\varphi_1^{-1}$ in $\varPsi$ und folglich $\varphi_2$ in $\varPsi\varphi_1$,
d.~h.\ in $\varPhi_1$ enthalten, gegen die Voraussetzung. Wir k"onnen
also eine ganze Zahl $\alpha$ in $\varOmega$ annehmen, die durch $\mathfrak{p} \mid \varphi_1$, aber
durch keines der Ideale $\mathfrak{p} \mid \varphi_2$ theilbar ist. Dann ist auch keine
der Zahlen $\alpha \mid \varphi'$ durch $\mathfrak{p} \mid \varphi_2$ theilbar, weil sonst $\mathfrak{p} \mid \varphi_1\varphi' = \mathfrak{p} \mid \varphi_2$,
also gegen die Voraussetzung $\varphi_2$ in $\varPhi_1$ enthalten w"are. Das
Product $\alpha'$ aller von einander verschiedener $\alpha \mid \varphi'$, welches eine
in $\varOmega'$ enthaltene Zahl ist, ist zwar durch $\mathfrak{p}_1'$, nicht aber durch $\mathfrak{p}_2'$
theilbar. Folglich ist $\mathfrak{p}_1'$ von $\mathfrak{p}_2'$ verschieden, und es ergiebt sich
also, dass $e'$ und nicht mehr \emph{verschiedene} Primideale $\mathfrak{p}'$ in
$\mathfrak{P}$ aufgehen. Wir setzen daher
\begin{equation}\label{eq:s163-3}
\mathfrak{P} = \mathfrak{p}_1'^{a_1}\mathfrak{p}_2'^{a_2}\dots\mathfrak{p}_{e'}^{a_{e'}},
\end{equation}
worin die Exponenten $a_1, a_2, \dots, a_{e'}$ noch n"aher zu bestimmende
nat"urliche Zahlen sind.

Um nun die Primideale $\mathfrak{p}_r'$ im K"orper $\varOmega$ zu zerlegen, k"onnen
wir die Resultate der \S.~160, 162 benutzen, indem wir einfach
$\varOmega'$ an Stelle von $R$ treten lassen. Bedeutet dann $X_r'$ den Durch-
schnitt von $\varPhi'$ mit $X_r = \varphi_r^{-1}X\varphi_r$, so ist $X_r'$ der Inbegriff aller
Substitutionen des K"orpers $\varPhi'$, die der Bedingung
\[
\tau \mid \chi_r' \equiv \tau \pmod{\mathfrak{p}_r}
\]
gen"ugen. Bezeichnen wir daher den Grad von $X_r'$ mit $g_r$, so ist
$\mathfrak{p}_r'$ durch $\mathfrak{p}^{g_r}$, aber durch keine h"ohere Potenz von $\mathfrak{p}$ theilbar.

Der Durchschnitt $\varPsi_r'$ von $\varPsi_r$ mit $\varPhi'$, dessen Grad wir schon
oben mit $h_r$ bezeichnet haben, ist der Inbegriff aller Substitu-
tionen $\psi_r'$ in $\varPhi'$, die der Bedingung
\[
\mathfrak{p}_r \mid \psi_r' = \mathfrak{p}_r
\]
gen"ugen, und wenn wir daher $\varPhi'$ in die Nebengruppen
\begin{equation}\label{eq:s163-4}
\varPhi' = \varPsi_r'\varphi_{r,1}' + \varPsi_r'\varphi_{r,2}' + \dots + \varPsi_r'\varphi_{r,e_r}'
\end{equation}
zerlegen, so ist
\begin{equation}\label{eq:s163-5}
n' = e_rh_r.
\end{equation}

Wenn nun
\[
\mathfrak{p}_{r,s} = \mathfrak{p}_r \mid \varphi_{r,s}
\]
gesetzt wird, so ist nach \S.~160,~(6)
\begin{equation}\label{eq:s163-6}
\mathfrak{p}_r' = (\mathfrak{p}_{r,1}\,\mathfrak{p}_{r,2}\dots\mathfrak{p}_{r,e_r})^{g_r},
\end{equation}
und die Substitution von~(6) in~(3) ergiebt, durch Vergleichung
mit \S.~160,~(6)
\begin{equation}\label{eq:s163-7}
g = a_rg_r, \quad e_1 + e_2 + \dots + e_{e'} = e,
\end{equation}
wodurch die Exponenten $a_r$ definirt sind.

Die in Bezug auf den K"orper $\varOmega'$ genommene Partialnorm
von $\mathfrak{p}_{r,s}$ ist nach \S.~160,~(3),~(7):
\begin{equation}\label{eq:s163-8}
\mathfrak{N}_{\varOmega'}(\mathfrak{p}_{r,s}) = \mathfrak{p}_r'^{f_r'},
\end{equation}
wenn $f_r$ durch die Gleichung
\begin{equation}\label{eq:s163-9}
n' = e_rf_rg_r, \quad (h_r = f_rg_r)
\end{equation}
definirt wird.

Wir wollen nun die Gruppe $\varPhi$ nach $\varPhi'$ in Nebengruppen
zerlegen, und setzen, wenn $\vartheta_1, \vartheta_2, \dots, \vartheta_{m'}$ passend ausgew"ahlte
Substitutionen aus $\varPhi$ sind:
\begin{equation}\label{eq:s163-10}
\varPhi = \varPhi'\vartheta_1 + \varPhi'\vartheta_2 + \dots + \varPhi'\vartheta_{m'},
\end{equation}
so dass durch die Substitutionen $\vartheta_1, \vartheta_2, \dots, \vartheta_{m'}$ der K"orper $\varOmega'$
in $m'$ conjugirte (gleiche oder verschiedene) K"orper
\[
\varOmega_1',\; \varOmega_2',\; \dots,\; \varOmega_{m'}'
\]
"ubergeht. Der K"orper $\varOmega_t'$ geh"ort dann zu der Gruppe $\vartheta_t^{-1}\varPhi'\vartheta_t$.

Wenn nun durch $\vartheta_t$ die Ideale $\mathfrak{p}_r'$ und $\mathfrak{p}_{r,s}$ in $\mathfrak{p}_{r,t}'$ und $\mathfrak{p}_{r,s,t}$
"ubergehen, so ist nach \S.~160,~(6)
\begin{equation}\label{eq:s163-11}
\mathfrak{p}_{r,t}' = (\mathfrak{p}_{r,1,t}\,\mathfrak{p}_{r,2,t}\dots\mathfrak{p}_{r,e_r,t})^{g_r},
\end{equation}
und das Product aller dieser Ideale f"ur $t = 1, 2, \dots, m'$ ist die
im K"orper $\varOmega'$ genommene Partialnorm von $\mathfrak{p}_r'$ in Bezug auf den
K"orper $R$, die wir mit $\mathfrak{N}_R'(\mathfrak{p}_r')$ bezeichnen. Sie muss eine Potenz
von $\mathfrak{P}$ sein, und wir setzen
\begin{equation}\label{eq:s163-12}
\mathfrak{N}_R(\mathfrak{p}_r') = \mathfrak{P}^{f_r'}.
\end{equation}

Um $f_r'$ zu finden, brauchen wir nur $\mathfrak{P}$ in~(12) nach \S.~160,~(6)
in seine Primfactoren $\mathfrak{p}$ zu zerlegen, wodurch sich $e\,g\,f_r'$ Prim-
factoren $\mathfrak{p}$ in $\mathfrak{P}^{f_r'}$ ergeben. Bilden wir andererseits das Product
 der $m'$ Factoren~(11), so erhalten wir $g_r\,e_r\,m'$ solcher Primfactoren.
Es ist daher
\[
e\,g\,f_r' = m'\,g_r\,e_r,
\]
also nach~(1),~(9) und \S.~160,~(7):
\begin{equation}\label{eq:s163-13}
f = f_r\,f_r'.
\end{equation}

\newpage
\section*{Zwanzigster Abschnitt.}
\addcontentsline{toc}{section}{Zwanzigster Abschnitt. Quadratische K"orper}

\begin{center}
\Large\bfseries
Quadratische K"orper.
\end{center}
\bigskip

\section*{\S.~164. Basis eines quadratischen K"orpers.}
\addcontentsline{toc}{section}{\S.~164. Basis eines quadratischen K"orpers}
\markboth{\S.~164. Basis eines quadratischen K"orpers}{\S.~164. Basis eines quadratischen K"orpers}

Wir wollen die Resultate der allgemeinen Theorie der alge-
braischen Zahlen auf den besonderen Fall der K"orper vom $2^{\mathrm{ten}}$
Grade, die auch \emph{quadratische K"orper} heissen, anwenden. Ein
quadratischer K"orper $\varOmega$ entspringt aus einer ganzzahligen
quadratischen Gleichung, und wird also aus dem K"orper der
rationalen Zahlen durch Adjunction einer Quadratwurzel $\sqrt{d}$ ab-
geleitet. Hierin kann $d$ als positive oder negative ganze Zahl
ohne quadratischen Theiler angenommen werden, und der ein-
zige Werth $d = 1$ ist auszunehmen. Je nachdem $d$ positiv oder
negativ ist, erhalten wir einen reellen oder imagin"aren quadra-
tischen K"orper.

Den conjugirten K"orper erhalten wir, wenn wir in allen
Zahlen in $\varOmega$ $\sqrt{d}$ in $-\sqrt{d}$ verwandeln. Dadurch bekommen wir
aber keinen neuen K"orper. Der quadratische K"orper ist also
ein (absoluter) Normalk"orper~(\S.~160).

Die Zahlen des K"orpers $\varOmega$ sind keine anderen als die in
\S.~122 des ersten Bandes definirten quadratischen Irrational-
zahlen, die alle in der Form $x + y\sqrt{d}$ enthalten sind, worin
$x$ und $y$ ganze oder gebrochene rationale Zahlen bedeuten. Jede
irrationale Zahl $\omega$ dieses K"orpers gen"ugt einer quadratischen
Gleichung
\begin{equation}\label{eq:s164-1}
c\,\omega^2 = a + b\,\omega,
\end{equation}
in der $a, b, c$ ganze rationale Zahlen ohne gemeinsamen Theiler
sind, von denen $c$ als positiv angenommen werden kann, und
die Verbindung
\begin{equation}\label{eq:s164-2}
D = b^2 + 4\,ac
\end{equation}
haben wir die \emph{Discriminante} der Zahl $\omega$ genannt. Das Ver-
h"altniss $D : d$ ist eine ganze rationale Quadratzahl. Die beiden
Wurzeln von~(1) erhalten den Ausdruck
\begin{equation}\label{eq:s164-3}
\omega = \frac{b + \sqrt{D}}{2c}, \quad \omega' = \frac{b - \sqrt{D}}{2c}.
\end{equation}

Setzen wir
\begin{equation}\label{eq:s164-4}
D = e^2d,
\end{equation}
so wird
\begin{equation}\label{eq:s164-5}
\omega = \frac{b + e\sqrt{d}}{2c}, \quad \omega' = \frac{b - e\sqrt{d}}{2c},
\end{equation}
und wegen~(1) sind dies dann und nur dann ganze Zahlen, wenn
$c = 1$ ist.

Nun ist $D \equiv 0$ oder $\equiv 1 \pmod{4}$, und wenn also $D$ gerade
ist, so ist es durch $4$ theilbar; $b$ ist in diesem Falle gerade,
und weil $d$ nicht durch $4$ theilbar sein kann, so muss auch $e$
gerade sein.

Ist aber $D$ ungerade, so sind auch $b$ und $e$ ungerade. Dieser
Fall kann aber, wie~(4) zeigt, nur dann eintreten, wenn $d \equiv 1
\pmod{4}$ ist.

Daraus erhalten wir folgendes Resultat:

\begin{satz}
Ist $d \equiv 2$ oder $\equiv 3 \pmod{4}$, so ist die Zahl
\[
\omega = x + y\sqrt{d}
\]
dann und nur dann eine ganze Zahl, wenn die
rationalen Zahlen $x, y$ ganze Zahlen sind. (Quadra-
tische Irrationalzahlen erster Art, nach Dirichlet.)
\end{satz}

\begin{satz}
Ist $d \equiv 1 \pmod{4}$, so ist
\[
\omega = \frac{x + y\sqrt{d}}{2}
\]
dann und nur dann eine ganze Zahl, wenn die
rationalen Zahlen $x, y$ entweder zwei gerade
oder zwei ungerade ganze Zahlen sind. (Quadra-
tische Irrationalzahlen erster oder zweiter Art, je nach-
dem $e$ gerade oder ungerade ist.)
\end{satz}

Letzteres k"onnen wir auch so ausdr"ucken:

\begin{satz}
Ist $d \equiv 1 \pmod{4}$, so ist
\[
\omega = x + y\,\frac{1+\sqrt{d}}{2}
\]
dann und nur dann eine ganze Zahl, wenn die
rationalen Zahlen $x$ und $y$ ganze Zahlen sind.
\end{satz}

\newpage
\section*{\S.~165. Die Primideale in $\varOmega$.}
\addcontentsline{toc}{section}{\S.~165. Die Primideale in $\varOmega$}
\markboth{\S.~165. Die Primideale in $\varOmega$}{\S.~165. Die Primideale in $\varOmega$}

Daraus ergiebt sich eine Basis von $\mathfrak{o}$:
\begin{alignat*}{2}
\text{a)}\quad & 1, \sqrt{d}, &\quad d &\equiv 2 \text{ oder } \equiv 3 \pmod{4}, \\
\text{b)}\quad & 1, \frac{1+\sqrt{d}}{2}, &\quad d &\equiv 1 \pmod{4}.
\end{alignat*}

Wenn wir also
\begin{equation}\label{eq:s164-6}
\begin{alignedat}{2}
\text{a)}\quad & \varTheta = \sqrt{d}, &\quad d &\equiv 2, 3 \pmod{4} \\
\text{b)}\quad & \varTheta = \frac{1+\sqrt{d}}{2}, &\quad d &\equiv 1 \pmod{4}
\end{alignedat}
\end{equation}
setzen, so ist $1, \varTheta$ eine Basis von~$\mathfrak{o}$.

F"ur die Grundzahl erhalten wir demnach, wenn $\varTheta'$ mit $\varTheta$
conjugirt ist,
\[
\varDelta = \begin{vmatrix} 1, & \varTheta \\ 1, & \varTheta' \end{vmatrix}^2 = (\varTheta' - \varTheta)^2,
\]
also
\begin{equation}\label{eq:s164-7}
\begin{alignedat}{2}
\text{a)}\quad & \varDelta = 4d, &\quad d &\equiv 2, 3 \pmod{4} \\
\text{b)}\quad & \varDelta = d, &\quad d &\equiv 1 \pmod{4},
\end{alignedat}
\end{equation}
\begin{equation}\label{eq:s164-8}
\varTheta - \varTheta' = \sqrt{\varDelta}.
\end{equation}

Die Grundzahl ist also $\equiv 0$ oder $\equiv 1 \pmod{4}$, und muss,
wenn sie durch $4$ theilbar ist, $\equiv 8, 12 \pmod{16}$ sein\footnote{Eine ganze rationale Zahl $\varDelta$, die nach dem Modul $4$ mit $0$ oder
mit $1$ congruent ist, hat Kronecker eine Zahl von Discriminantenform
genannt. Hat $\varDelta$ die Eigenschaft, dass sich kein quadratischer Factor so
daraus absondern l"asst, dass eine Zahl von Discriminantenform "ubrig
bleibt, so heisst $\varDelta$ eine \emph{Stammdiscriminante}. Die Grundzahl eines
quadratischen K"orpers ist also immer eine Stammdiscriminante, und
umgekehrt ist auch jede Stammdiscriminante Grundzahl eines quadra-
tischen K"orpers. Vergl.\ H.~Weber, \emph{Zahlentheoretische Untersuchungen
aus dem Gebiete der elliptischen Functionen.} I, II, III. Nachrichten der
Gesellschaft der Wissenschaften zu G"ottingen. 1893.}. Sie kann
daher (in Uebereinstimmung mit dem Satze von Minkowski)
weder $= \pm 1$ noch $= \pm 2$ sein. Die Grundzahl vom absolut
kleinsten Werthe ist $\varDelta = -3$.

\bigskip

Es bedeute nun $\varphi$ ein beliebiges ganzes Functional in~$\varOmega$.
Wir wollen nach \S.~146 eine Basis $\alpha_1, \alpha_2$ f"ur $\varphi$, und damit die
entsprechende Basisform $\alpha_1t_1 + \alpha_2t_2$, die mit $\varphi$ associirt ist, be-
stimmen. Wir haben nach \S.~146,~(3), wenn $\omega_1, \omega_2$ eine Basis
von $\mathfrak{o}$ ist, die ganzen Zahlen $a_{1,1}, a_{1,2}, a_{2,2}$ so zu bestimmen, dass
$a_{1,1}, a_{2,2}$ positiv und m"oglichst klein sind, und dass
\[
\alpha_1 = a_{1,1}\omega_1, \quad \alpha_2 = a_{1,2}\omega_1 + a_{2,2}\omega_2
\]
durch $\varphi$ theilbar werden. Setzen wir also $(\omega_1, \omega_2) = (1, \varTheta)$, so
wird
\begin{equation}\label{eq:s165-1}
\alpha_1 = a_{1,1}, \quad \alpha_2 = a_{1,2} + a_{2,2}\varTheta.
\end{equation}
Es ist also $a_{1,1}$ die kleinste positive durch $\varphi$ theilbare rationale
Zahl, und $a_{2,2}$ ist die kleinste positive Zahl, f"ur die $a_{2,2}\varTheta$ nach
dem Modul $\varphi$ mit einer rationalen Zahl congruent wird.

Wir wollen dies auf den Fall anwenden, dass $\varphi$ ein Prim-
functional $\pi$ ist. Ein solches kann hier nur vom ersten oder
vom zweiten Grade sein, je nachdem seine absolute Norm gleich
$p$ oder $= p^2$ ist. Wir unterscheiden also die beiden F"alle:
\begin{alignat*}{2}
\text{1)}\quad & N_a(\pi) = p & \quad \text{oder} \\
\text{2)}\quad & N_a(\pi) = p^2. &
\end{alignat*}

Da hier $N(p) = p^2$ ist, so ist im ersten Falle $p$ in zwei
Primfactoren zerlegbar, im zweiten Falle enth"alt $p$ nur einen
Primfactor, d.~h.\ $p$ ist im K"orper $\varOmega$ selbst noch als Primzahl zu
betrachten.

Bestimmen wir nun f"ur $\varphi = \pi$ die Basis~(1), so ergiebt
sich zun"achst $a_{1,1} = p$, und $a_{2,2}$ ist die kleinste positive ganze
rationale Zahl, f"ur die das Product $a_{2,2}\varTheta$ modulo $\pi$ mit einer
rationalen Zahl congruent wird. Daf"ur ist aber nach \S.~150,~3
die nothwendige und hinreichende Bedingung
\[
(a_{2,2}\varTheta)^p \equiv a_{2,2}\varTheta \pmod{\pi},
\]
oder, da nach dem Fermat'schen Satze $a_{2,2}^p \equiv a_{2,2}$ ist,
\begin{equation}\label{eq:s165-2}
a_{2,2}(\varTheta^p - \varTheta) \equiv 0 \pmod{\pi}.
\end{equation}

Es ist also entweder $\varTheta^p - \varTheta \equiv 0$, und dann ist $a_{2,2} = 1$,
oder $\varTheta^p - \varTheta$ nicht $\equiv 0$, und dann ist $a_{2,2} = p$.

Im ersten Falle ist $\varTheta$ und damit jede Zahl $\omega$ mit einer
rationalen Zahl congruent. Es giebt also nicht mehr als $p$ nach
dem Modul $\pi$ incongruente Zahlen, und folglich ist~(\S.~148,~3.)
$p^f = p$, $f = 1$, d.~h.\ $p$ in zwei Primfactoren ersten Grades
zerlegbar.

\newpage

Umgekehrt bilden in diesem Falle~1) die $p$ rationalen Zahlen
$0, 1, 2, \dots, p-1$ ein volles Restsystem nach dem Modul $\pi$, und
es ist folglich jede Zahl in $\mathfrak{o}$ und mithin auch $\varTheta$ mit einer ratio-
nalen Zahl congruent. Also haben wir
\begin{alignat*}{2}
\text{1)}\quad & N_a(\pi) = p & \quad a_{2,2} &= 1 \\
\text{2)}\quad & N_a(\pi) = p^2 & \quad a_{2,2} &= p.
\end{alignat*}

Im ersten Falle ist $a_{1,2}$ aus der Bedingung
\begin{equation}\label{eq:s165-3}
a_{1,2} + \varTheta \equiv 0 \pmod{\pi}
\end{equation}
zu bestimmen, im zweiten kann $a_{1,2} = 0$ genommen werden. Die
Basisform von $\pi$ ergiebt sich also, wenn $u, v$ die Variablen sind,
\begin{alignat*}{2}
\text{1)}\quad & \pi = p\,u + (a_{1,2} + \varTheta)\,v \\
\text{2)}\quad & \pi = p\,(u + \varTheta\,v),
\end{alignat*}
woraus man sieht, dass $\pi$ im Falle~2), wie zu erwarten war, ein
\emph{Hauptfunctional} ist.

Um im Falle~1) die Zahl $a_{1,2}$ zu bestimmen, bilden wir die
Norm von $a_{1,2} + \varTheta$, die durch $p$ theilbar sein muss. Wenn $\varTheta$
durch Aenderung des Vorzeichens von $\sqrt{d}$ in $\varTheta'$ "ubergeht, so ist
nach \S.~164,~(8) $\varTheta - \varTheta' = \sqrt{\varDelta}$, und f"ur die Norm von $a_{1,2} + \varTheta$
ergiebt sich nach \S.~164~(6), a), b):
\begin{equation}\label{eq:s165-4}
\begin{aligned}
N(a_{1,2} + \varTheta) &= (a_{1,2} + \varTheta)(a_{1,2} + \varTheta') \\
&= a_{1,2}^2 - d, \quad d \equiv 2, 3 \pmod{4} \\
&= a_{1,2}^2 + a_{1,2} + \frac{1-d}{4}, \quad d \equiv 1 \pmod{4}.
\end{aligned}
\end{equation}

Wenn diese Zahl durch $p$ theilbar ist, so ist einer der
beiden Factoren $a_{1,2} + \varTheta$, $a_{1,2} + \varTheta'$ durch $\pi$ theilbar. Da aber
$\varTheta' = -\varTheta$ oder $= 1 - \varTheta$ ist, also
\begin{alignat*}{2}
\text{a)}\quad & a_{1,2} + \varTheta' = -(-a_{1,2} + \varTheta) \\
\text{b)}\quad & a_{1,2} + \varTheta' = -(-a_{1,2} - 1 + \varTheta)
\end{alignat*}
ist, so kann man $a_{1,2}$ immer so annehmen, dass gerade $a_{1,2} + \varTheta$
durch $\pi$ theilbar wird, und $\varTheta$ ist also wirklich mit einer ratio-
nalen Zahl congruent. Die nothwendige und hinreichende Be-
dingung daf"ur, dass $p$ zwei Primfactoren enth"alt, ist folglich die,
dass $a_{1,2}^2 - d$ oder $a_{1,2}^2 + a_{1,2} + \frac{1}{4}(1-d)$ durch geeignete An-
nahme "uber $a_{1,2}$ durch $p$ theilbar wird. Durch Multiplication
mit $4$ gehen diese beiden Bedingungen "uber in
\[
4a_{1,2}^2 - 4d \equiv 0, \quad (2a_{1,2} + 1)^2 - d \equiv 0 \pmod{4p},
\]

und wenn man bedenkt, dass die Grundzahl $\varDelta$ im Falle a)
gleich $4d$, im Falle b) gleich $d$ ist, so erhalten wir, wenn wir
\begin{equation}\label{eq:s165-5}
x = 2a_{1,2} \quad \text{oder} \quad = 2a_{1,2} + 1
\end{equation}
setzen, folgendes Resultat:

\begin{satz}
Die nothwendige und hinreichende Bedingung
daf"ur, dass $p$ zwei Primfactoren enth"alt, ist
die, dass die Congruenz
\begin{equation}\label{eq:s165-6}
x^2 \equiv \varDelta \pmod{4p}
\end{equation}
durch eine ganze rationale Zahl $x$ befriedigt
werden kann, oder dass die Grundzahl $\varDelta$ qua-
dratischer Rest von $4p$ ist.
\end{satz}

Bezeichnen wir mit $\pi, \pi'$ die beiden conjugirten Formen
\begin{equation}\label{eq:s165-7}
\pi = p\,u + (a_{1,2} + \varTheta)\,v, \quad \pi' = p\,u + (a_{1,2} + \varTheta')\,v,
\end{equation}
so ergiebt sich durch Multiplication
\begin{equation}\label{eq:s165-8}
\pi\pi' = p\left(p\,u^2 + x\,u\,v + \frac{x^2 - \varDelta}{4p}\,v^2\right).
\end{equation}
Hierin ist nun
\[
p\,u^2 + x\,u\,v + \frac{x^2 - \varDelta}{4p}\,v^2
\]
eine Einheit, weil, wenn $p$ in $\varDelta$ nicht aufgeht, $x$ durch $p$ nicht
theilbar ist, und wenn $p$ in $\varDelta$ aufgeht, zwar $x$, aber nicht
$(x^2 - \varDelta : 4p)$ durch $p$ theilbar ist. Denn ist $p$ ungerade, so kann
$p^2$ nicht in $\varDelta$ aufgehen, und ist $p = 2$ und $\varDelta$ durch $2$ theilbar,
so ist $\varDelta = 4d$, $d \equiv 2, 3 \pmod{4}$. W"are aber
\[
x^2 - \varDelta \equiv 0 \pmod{16},
\]
so w"urde folgen
\[
d \equiv \left(\frac{x}{2}\right)^2 \pmod{4},
\]
w"ahrend doch $\left(\dfrac{x}{2}\right)^2$ als Quadrat nicht mit $2$ oder $3$ congruent
sein kann. Demnach ist $p$ in die beiden Primfactoren $\pi$ und $\pi'$
zerlegt.

Die beiden Factoren $\pi$ und $\pi'$ von $p$ k"onnen auch mit
einander associirt sein. Dies tritt nur dann ein, wenn $\pi'$ und
folglich auch $\pi - \pi'$ durch $\pi$ theilbar ist. Nun ist aber nach~(7)
\[
\pi - \pi' = (\varTheta - \varTheta')\,v,
\]
und folglich muss
\[
(\varTheta - \varTheta')^2 = \varDelta
\]
durch $p$ theilbar sein, und dies ist auch die hinreichende Be-
dingung.

\begin{satz}
Es ist also $p$ nur dann mit dem Quadrat eines
Primfunctionals associirt, wenn $p$ in der Grundzahl
$\varDelta$ aufgeht, in Uebereinstimmung mit dem
allgemeinen Satze~(\S.~158).
\end{satz}

\newpage
\section*{\S.~166. Functionale im quadratischen K"orper und quadratische Irrationalzahlen.}
\addcontentsline{toc}{section}{\S.~166. Functionale im quadratischen K"orper und quadratische Irrationalzahlen}
\markboth{\S.~166. Functionale im quadratischen K"orper}{\S.~166. Functionale im quadratischen K"orper}

Wir betrachten nun in unserem quadratischen K"orper ein
beliebiges ganzes Functional und bilden seine Basisform
\[
\varphi = \alpha_1t_1 + \alpha_2t_2,
\]
indem wir $\alpha_1, \alpha_2$ nach \S.~165 bestimmen:
\begin{equation}\label{eq:s166-1}
\alpha_1 = a_{1,1}, \quad \alpha_2 = a_{1,2} + a_{2,2}\varTheta.
\end{equation}

Nach \S.~146 ist dann
\begin{equation}\label{eq:s166-2}
N_a(\varphi) = a_{1,1}\,a_{2,2}.
\end{equation}

Geben wir der Form $\varphi$ den Ausdruck
\begin{equation}\label{eq:s166-3}
\varphi = a_{1,1}t_1 + (a_{1,2} + a_{2,2}\varTheta)\,t_2,
\end{equation}
und bilden die Norm, so ergiebt sich
\begin{equation}\label{eq:s166-4}
N(\varphi) = a_{1,1}^2t_1^2 + [2a_{1,1}a_{1,2} + a_{1,1}a_{2,2}(\varTheta + \varTheta')]\,t_1t_2 + t_2^2\,N(\alpha_2)
\end{equation}
und der Theiler dieser Form ist also nach~(2) gleich $a_{1,1}a_{2,2}$.

Wenn wir also
\begin{equation}\label{eq:s166-5}
\begin{aligned}
a_{1,1}^2 &= c\;a_{1,1}\,a_{2,2}, \\
2a_{1,1}a_{1,2} + a_{1,1}a_{2,2}(\varTheta + \varTheta') &= b\;a_{1,1}\,a_{2,2}, \\
N(\alpha_2) &= -a\;a_{1,1}\,a_{2,2}
\end{aligned}
\end{equation}
setzen, so sind $a, b, c$ ganze rationale Zahlen ohne gemeinschaft-
lichen Theiler, und der Quotient
\begin{equation}\label{eq:s166-6}
\frac{\alpha_2}{\alpha_1} = \omega
\end{equation}
ist eine Wurzel der quadratischen Gleichung
\begin{equation}\label{eq:s166-7}
c\,\omega^2 = a + b\,\omega.
\end{equation}
Wenn wir in den Ausdr"ucken~(1) die Werthe von $a_{1,1}$ und $a_{1,2}$
aus~(5) substituiren, so folgt
\begin{equation}\label{eq:s166-8}
\begin{aligned}
2\alpha_2 &= a_{2,2}\,(b + \varTheta - \varTheta'), \\
\alpha_1 &= c\,a_{2,2}.
\end{aligned}
\end{equation}

Nun ist nach \S.~164,~(8) die Differenz $\varTheta - \varTheta'$ gleich der
Quadratwurzel aus der Grundzahl von $\varOmega$, so dass wir aus~(6)
und~(8) erhalten:
\begin{equation}\label{eq:s166-9}
\omega = \frac{b + \sqrt{\varDelta}}{2c}, \quad \varDelta = b^2 + 4\,ac,
\end{equation}
worin das Vorzeichen von $\sqrt{\varDelta}$ durch die in \S.~164,~(6) gemachte
Annahme "uber das Vorzeichen von $\sqrt{d}$ bestimmt ist.

Es ist also $\omega$ eine zur Discriminante $\varDelta$ geh"orige quadra-
tische Irrationalzahl, wie wir sie im \S.~122 des ersten Bandes
betrachtet haben, und die Form $\varphi$ kann, wenn wir $a_{2,2} = e$
setzen, so ausgedr"uckt werden:
\begin{equation}\label{eq:s166-10}
\varphi = e\,c\,(t_1 + \omega\,t_2).
\end{equation}

Hierin ist $c$ die kleinste positive ganze rationale Zahl, f"ur
die das Product $c\,\omega$ eine ganze Zahl wird, und jede andere
ganze rationale Zahl $k$, durch die das Product $k\,\omega$ eine ganze
Zahl wird, muss durch $c$ theilbar sein.

Man erkennt dies am einfachsten, wenn man $\omega$ nach \S.~164,
je nachdem $b$ gerade oder ungerade ist, in eine der beiden
Formen setzt:
\begin{equation}\label{eq:s166-11}
\omega = \frac{1}{c}\left(\frac{b}{2} + \varTheta\right), \quad \frac{1}{c}\left(\frac{b-1}{2} + \varTheta\right),
\end{equation}
woraus, da $1, \varTheta$ eine Basis von $\mathfrak{o}$ ist, das Gesagte hervorgeht.

Darauf gr"undet sich nun der Beweis, dass, wenn $\omega$ irgend
eine zur Discriminante $\varDelta$ geh"orige quadratische Irrationalzahl
bedeutet, $c$ die kleinste positive ganze rationale Zahl, die $c\,\omega$ zu
einer ganzen Zahl macht, und $e$ eine beliebige positive ganze
rationale Zahl, auch umgekehrt durch
\begin{equation}\label{eq:s166-12}
\varphi = e\,c\,(t_1 + \omega\,t_2)
\end{equation}
immer eine Basisform dargestellt ist.

Dazu haben wir nur zu zeigen, dass
\[
\alpha_1 = e\,c, \quad \alpha_2 = e\,c\,\omega
\]
eine Basis des durch~(12) definirten Functionals $\varphi$ ist.

Bilden wir zu diesem Zwecke zun"achst die Norm von $\varphi$, so
erhalten wir nach~(7)
\begin{equation}\label{eq:s166-13}
N(\varphi) = e^2c\,(c\,t_1^2 + b\,t_1t_2 - a\,t_2^2),
\end{equation}
und folglich die absolute Norm
\begin{equation}\label{eq:s166-14}
N_a(\varphi) = e^2c.
\end{equation}

Ist nun
\begin{equation}\label{eq:s166-15}
\alpha_1 = a_{1,1}, \quad \alpha_2 = a_{1,2} + a_{2,2}\varTheta
\end{equation}
eine Basis von $\varphi$, so ist andererseits (nach \S.~146)
\begin{equation}\label{eq:s166-16}
N_a(\varphi) = a_{1,1}\,a_{2,2},
\end{equation}
und $a_{1,1}$ ist die kleinste positive ganze rationale Zahl, die durch
$\varphi$ theilbar ist.

Nun k"onnen wir andererseits zeigen, dass diese Zahl $= e\,c$
ist, woraus $e\,c = a_{1,1}$, und nach~(14) und~(16) $e = a_{2,2}$ folgt.

Wenn n"amlich eine ganze rationale Zahl $m$ durch $\varphi$ theil-
bar sein soll, so muss, wenn $\omega'$ zu $\omega$ conjugirt ist, da
\[
c\,t_1^2 + b\,t_1t_2 - a\,t_2^2
\]
ein Einheitsfunctional ist,
\[
\frac{m\,(c\,t_1^2 + b\,t_1t_2 - a\,t_2^2)}{\varphi} = \frac{m}{e}\,(t_1 + \omega'\,t_2)
\]
ein ganzes Functional sein. Folglich m"ussen
\[
\frac{m}{e} \quad \text{und} \quad \frac{m}{e}\,\omega'
\]
ganze Zahlen sein, woraus folgt, dass $m : e$ eine durch $c$ theil-
bare, also $m$ eine durch $e\,c$ theilbare ganze rationale Zahl ist.
Da andererseits $e\,c$ durch $\varphi$ theilbar ist, so folgt, dass $e\,c$ die
\emph{kleinste} diesen Forderungen gen"ugende Zahl ist. Ausserdem er-
giebt sich noch, dass die ganze Zahl $e\,c\,\omega$ durch $\varphi$ theilbar ist.
Denn $e\,c\,\omega$ ist in der Darstellung~(12) ein Coefficient der ganzen
Function $\varphi$ und daher durch $\varphi$ theilbar~(\S.~142).

Dr"ucken wir $\varTheta$ in~(15) durch $\omega$ aus, so ergiebt sich nach~(11)
\begin{alignat*}{2}
\alpha_2 &= a_{1,2} - \frac{eb}{2} + e\,c\,\omega, &\quad &b \text{ gerade} \\
         &= a_{1,2} - \frac{e(b-1)}{2} + e\,c\,\omega, &\quad &b \text{ ungerade},
\end{alignat*}
und es ist also
\[
a_{1,2} \equiv \frac{eb}{2} \quad \text{oder} \quad \equiv \frac{e(b-1)}{2}
\]
nach dem Modul $\varphi$, und folglich, da beide Seiten rational sind,
auch nach dem Modul $a_{1,1} = e\,c$. Demnach kann $a_{1,2}$ dem
Werthe $\frac{1}{2}eb$ oder $\frac{1}{2}e(b-1)$ gleich gesetzt werden, und wir er-
halten $\alpha_2 = e\,c\,\omega$, was zu beweisen war.

Um das hierdurch Beweisene "ubersichtlich zusammenzu-
fassen, k"onnen wir folgende S"atze aussprechen:

\begin{satz}
Ist $\varOmega$ ein quadratischer K"orper mit der Grund-
zahl $\varDelta$, und $\varphi$ ein ganzes Functional in diesem
K"orper, so l"asst sich eine zur Discriminante $\varDelta$
geh"orige quadratische Irrationalzahl
\begin{equation}\label{eq:s166-17}
\omega = \frac{b + \sqrt{\varDelta}}{2c}
\end{equation}
und eine positive ganze rationale Zahl $e$ so
bestimmen, dass
\begin{equation}\label{eq:s166-18}
e\,c\,(t_1 + \omega\,t_2)
\end{equation}
eine Basisform von $\varphi$ wird. Umgekehrt ist, wenn
$\omega$ eine beliebige zur Discriminante $\varDelta$ geh"orige
quadratische Irrationalzahl ist, die Form~(18)
immer eine Basisform eines Functionals~$\varphi$.
\end{satz}

\newpage
\section*{\S.~167. Aequivalente Functionale und "aquivalente Zahlen.}
\addcontentsline{toc}{section}{\S.~167. Aequivalente Functionale und "aquivalente Zahlen}
\markboth{\S.~167. Aequivalente Functionale und "aquivalente Zahlen}{\S.~167. Aequivalente Functionale und "aquivalente Zahlen}

Wenn nun $\omega, \omega_1$ irgend zwei zu der Discriminante $\varDelta$ geh"orige
quadratische Irrationalzahlen sind, so erhalten wir zwei Basis-
formen
\begin{equation}\label{eq:s167-1}
\begin{aligned}
\varphi  &= e\,c\,(t_1 + \omega\,t_2), \\
\varphi_1 &= e_1c_1\,(t_1 + \omega_1t_2).
\end{aligned}
\end{equation}

Wenn nun die beiden Zahlen $\omega, \omega_1$ in dem Sinne "aqui-
valent sind, wie wir diesen Begriff im \S.~120 des ersten Bandes
festgesetzt haben, so ist
\begin{equation}\label{eq:s167-2}
\omega_1 = \frac{p\,\omega + q}{r\,\omega + s},
\end{equation}
worin $p, q, r, s$ ganze rationale Zahlen sind, die der Bedingung
\begin{equation}\label{eq:s167-3}
ps - qr = \pm 1
\end{equation}
gen"ugen. Machen wir dann die Substitution~(2) in $\varphi_1$, so folgt
\[
(r\omega + s)\,\varphi_1 = e_1c_1\,[s\,t_1 + q\,t_2 + \omega\,(r\,t_1 + p\,t_2)],
\]
oder wenn wir
\begin{equation}\label{eq:s167-4}
u_1 = s\,t_1 + q\,t_2, \quad u_2 = r\,t_1 + p\,t_2
\end{equation}
setzen,
\[
(r\omega + s)\,\varphi_1 = e_1c_1\,(u_1 + \omega\,u_2).
\]

Nun ist $e\,c\,(u_1 + \omega\,u_2)$ mit $\varphi$ associirt, weil~(4) eine ganz-
zahlige lineare Substitution mit der Determinante $\pm 1$ ist~(\S.~147),
und folglich sind die beiden Formen
\[
(r\omega + s)\,e\,c\,\varphi_1, \quad e_1c_1\,\varphi
\]
mit einander associirt, und dies hat zur Folge, dass die beiden
Formen $\varphi, \varphi_1$ in dem in \S.~152 definirten Sinne "aquivalent sind.
Damit haben wir bewiesen:

\begin{satz}
Wenn die quadratischen Irrationalzahlen $\omega, \omega_1$
"aquivalent sind, so sind auch die entsprechenden
Functionale $\varphi, \varphi_1$ "aquivalent.
\end{satz}

Es l"asst sich aber auch der umgekehrte Satz beweisen:

\begin{satz}
Wenn die Functionale $\varphi, \varphi_1$ "aquivalent sind, so
sind auch die quadratischen Irrationalzahlen
$\omega$ und $\omega_1$ "aquivalent.
\end{satz}

Um diesen Beweis zu f"uhren, bezeichnen wir mit $\psi$ ein
Functional der Classe, die zu der Classe von $\varphi$ reciprok ist,
so dass $\varphi\,\psi$ mit einer Zahl "aquivalent ist. Dann k"onnen wir
setzen:
\begin{equation}\label{eq:s167-5}
\varphi\,\psi = \varepsilon\,\xi,
\end{equation}
worin $\varepsilon$ eine Einheit, $\xi$ eine Zahl in $\mathfrak{o}$ ist. Hierin k"onnen wir
$\varphi$ und $\psi$ beide als Basisformen annehmen und demnach
\begin{equation}\label{eq:s167-6}
\varphi = e\,c\,(t_1 + t_2\,\omega), \quad \psi = e'c'\,(t_1 + t_2\,\eta)
\end{equation}
setzen, worin $e', c', \eta$ dieselbe Bedeutung f"ur $\psi$ haben, wie $e, c, \omega$
f"ur $\varphi$. Wir setzen zur Abk"urzung [\S.~166,~(13)]:
\begin{equation}\label{eq:s167-7}
N_a(\varphi) = e^2c = P, \quad N_a(\psi) = e'^2c' = Q.
\end{equation}

Wenn wir mit $\psi'$ die zu $\psi$ conjugirte Form bezeichnen, und
f"ur den Augenblick die Einheitsform
\[
c't_1^2 + b't_1t_2 - a't_2^2 = \varepsilon_1
\]
setzen, so folgt aus~(5) durch Multiplication mit $\psi'$:
\[
Q\,\varepsilon_1\,\varphi = \varepsilon\,\xi\,\psi',
\]
und folglich
\begin{equation}\label{eq:s167-8}
\frac{Q\,\varphi}{\xi} = \frac{\varepsilon}{\varepsilon_1}\,\psi' = \chi,
\end{equation}
worin $\chi$ und $\psi'$ associirt sind. Setzen wir
\begin{equation}\label{eq:s167-9}
\frac{Q\,e\,c}{\xi} = \beta_1, \quad \frac{Q\,e\,c\,\omega}{\xi} = \beta_2,
\end{equation}
so wird
\[
\chi = \frac{Q\,\varphi}{\xi} = \beta_1t_1 + \beta_2t_2,
\]
und folglich sind $\beta_1, \beta_2$ ganze durch $\psi'$ theilbare Zahlen.

Bilden wir aus~(9) die Discriminante $\varDelta(\beta_1, \beta_2)$, so er-
halten wir
\begin{equation}\label{eq:s167-10}
\varDelta(\beta_1, \beta_2) = \begin{vmatrix} \beta_1, & \beta_2 \\ \beta_1', & \beta_2' \end{vmatrix}^2 = \frac{e^4c^4\,Q^4}{(\xi\xi')^2}\,(\omega - \omega')^2.
\end{equation}
Aus~(5) und~(7) folgt aber
\[
\xi\xi' = N(\xi) = PQ = e^2c\,Q,
\]
und wenn wir noch $\omega - \omega' = \dfrac{\sqrt{\varDelta}}{c}$ setzen, so findet sich
\begin{equation}\label{eq:s167-11}
\varDelta(\beta_1, \beta_2) = Q^2\varDelta.
\end{equation}

Hieraus folgt aber nach dem Satze \S.~147, 2., dass $\beta_1, \beta_2$
eine Basis von $\psi'$ ist.

Andererseits ist nach~(6) auch $e'c', e'c'\eta'$ eine Basis von $\psi'$,
und folglich lassen sich die ganzen rationalen Zahlen $p, q, r, s$
so bestimmen, dass
\begin{equation}\label{eq:s167-12}
\begin{aligned}
\beta_1 = \frac{Q\,e\,c}{\xi} &= e'c'\,(r\eta' + s), \\
\beta_2 = \frac{Q\,e\,c}{\xi}\,\omega &= e'c'\,(p\eta' + q),
\end{aligned}
\end{equation}
und zugleich
\begin{equation}\label{eq:s167-13}
ps - qr = \pm 1.
\end{equation}
Aus~(12) folgt aber durch Division:
\begin{equation}\label{eq:s167-14}
\omega = \frac{p\eta' + q}{r\eta' + s},
\end{equation}
d.~h.\ $\omega$ ist mit $\eta'$ "aquivalent.

Ersetzen wir nun $\varphi$ durch die "aquivalente Form $\varphi_1$, so
besteht eine Gleichung wie~(5):
\[
\varphi_1\,\psi = \varepsilon_1\xi_1,
\]
worin $\psi$ unge"andert geblieben ist. Folglich ist auch $\omega_1$ mit $\eta$
"aquivalent, und mithin sind $\omega$ und $\omega_1$ unter einander "aqui-
valent. Das ist aber die in unserem Satze~3.\ ausgesprochene
Behauptung, die also hiermit erwiesen ist.

Hiernach entspricht also jeder Idealclasse des K"orpers $\varOmega$
eine Classe "aquivalenter Zahlen $\omega$ und umgekehrt, und wir
k"onnen daher noch den Satz hinzuf"ugen:

\begin{satz}
Die Anzahl der nicht "aquivalenten Irrational-
zahlen der Discriminante $\varDelta$ ist gleich der Anzahl
der Formenclasses des K"orpers~$\varOmega$\footnote{Nach der am Schlusse des \S.~128 des ersten Bandes gemachten Be-
merkung "uber die Gauss'sche Theorie der quadratischen Formen stimmt
die Classenzahl der quadratischen Formen nach der Gauss'schen Defini-
tion nicht immer "uberein mit der Anzahl der Idealclasses des K"orpers $\varOmega$.
Um die Uebereinstimmung herzustellen, muss man in manchen F"allen die
Idealclasses noch weiter in je zwei Classen theilen. Bei unserer Theorie
der Aequivalenz der Zahlen $\omega$, wo die eigentliche nicht von der uneigent-
lichen Aequivalenz unterschieden wird, ist diese weitere Eintheilung nicht
erforderlich. Vergl.\ Dirichlet-Dedekind, Zahlentheorie. 4.~Auflage,
S.~578, 584.}.
\end{satz}

\newpage
\section*{\S.~168. Composition der quadratischen Irrationalzahlen.}
\addcontentsline{toc}{section}{\S.~168. Composition der quadratischen Irrationalzahlen}
\markboth{\S.~168. Composition der quadratischen Irrationalzahlen}{\S.~168. Composition der quadratischen Irrationalzahlen}

Wir haben im \S.~166 gesehen, dass zu jedem ganzen Func-
tional $\varphi$ im K"orper $\varOmega$ eine gewisse quadratische Irrationalzahl $\omega$
von der Discriminante $\varDelta$ gefunden werden kann.

Wir m"ussen aber jetzt noch untersuchen, inwiefern diese
Zahl $\omega$ durch den K"orper $\varOmega$ und durch das Functional $\varphi$ be-
stimmt ist.

Nach \S.~166,~(1) ist die Basis von $\varphi$
\begin{equation}\label{eq:s168-1}
\alpha_1 = a_{1,1}, \; \alpha_2 = a_{1,2} + a_{2,2}\varTheta
\end{equation}
dadurch definirt, dass $a_{1,1}$ die kleinste positive ganze rationale
Zahl ist, die durch $\varphi$ theilbar ist, und $a_{2,2}$ ist dann durch
\begin{equation}\label{eq:s168-2}
N_a(\varphi) = a_{1,1}\,a_{2,2}
\end{equation}
bestimmt, also sind die beiden Zahlen $a_{1,1}, a_{2,2}$ durch $\varphi$ ein-
deutig bestimmt. Es fragt sich noch, inwiefern auch $a_{1,2}$ be-
stimmt ist. Die Zahl $a_{1,2}$ ist aber nach dem Modul $\varphi$ durch die
Congruenz $a_{1,2} + a_{2,2}\varTheta \equiv 0$ gegeben, und da $a_{1,2}$ auch eine
ganze rationale Zahl sein muss, so ist $a_{1,2}$ bis auf ein Vielfaches
von $a_{1,1}$ bestimmt. Dieses Vielfache von $a_{1,1}$ bleibt aber in der
That willk"urlich, da, wenn $\alpha_1, \alpha_2$ eine Basis von $\varphi$ ist, dasselbe
f"ur jedes ganze rationale $k$ auch von $\alpha_1, \alpha_2 + k\alpha_1$ gilt.

Da nun $\alpha_2 : \alpha_1 = \omega$ die gesuchte quadratische Irrational-
zahl ist, so erhalten wir den Satz:

\begin{satz}
Zu jedem ganzen Functional des K"orpers $\varOmega$ geh"ort
eine und nur eine Reihe von quadratischen
Irrationalzahlen der Discriminante $\varDelta$, und die
Zahlen dieser Reihe unterscheiden sich um
beliebige ganze rationale Zahlen und sind
folglich unter einander "aquivalent. Associirte
Formen f"uhren auf dieselbe Reihe.
\end{satz}

Wenn $\varphi_1, \varphi_2$ zwei ganze Functionale in $\varOmega$ bedeuten, und
\[
\varphi_1\,\varphi_2 = \varphi_3
\]
gesetzt wird, so k"onnen wir zu jeder dieser drei Formen eine
quadratische Irrationalzahl $\omega_1, \omega_2, \omega_3$ bestimmen, und $\omega_3$ ist
durch $\omega_1, \omega_2$ bis auf eine additive ganze rationale Zahl bestimmt.
Man nennt dann $\omega_3$ aus $\omega_1$ und $\omega_2$ \emph{componirt} oder zusammen-
gesetzt, und setzt in symbolischer Bezeichnung
\begin{equation}\label{eq:s168-3}
\omega_1\,\omega_2 = \omega_3.
\end{equation}

Ist $\psi_1$ "aquivalent mit $\varphi_1$, $\psi_2$ mit $\varphi_2$, so ist nach \S.~152 auch
$\psi_1\psi_2 = \psi_3$ mit $\varphi_3$ "aquivalent.

Wenn nun $\eta_1, \eta_2, \eta_3$ die durch den Satz~5.\ zu den Formen
$\psi_1, \psi_2, \psi_3$ geh"origen Irrationalzahlen sind, so bilden nach~3.,
\S.~167 die Zahlen $\eta_1, \omega_1$; $\eta_2, \omega_2$; $\eta_3, \omega_3$ drei Paare "aquivalenter
Zahlen und wir k"onnen daher folgenden Satz aussprechen:

\begin{satz}
Ersetzt man in einer Composition von quadra-
tischen Irrationalzahlen jedes Element durch
eine "aquivalente Zahl, so geht auch die com-
ponirte Zahl in eine "aquivalente "uber.
\end{satz}

Dadurch sind die Classen der quadratischen Irrationalzahlen
einer Discriminante $\varDelta$ zu einer mit der Gruppe der Idealclassen
isomorphen Abel'schen Gruppe verbunden.

Um die aus zwei gegebenen Zahlen componirte Zahl wirk-
lich zu finden, hat man so zu verfahren.

Wir wollen die nach~(1) gebildeten Basen der Formen
$\varphi_1, \varphi_2, \varphi_3$ so bezeichnen:
\[
\begin{array}{lll}
\varphi_1)\quad & \alpha_1 = a_{1,1}, & \alpha_2 = a_{1,2} + a_{2,2}\varTheta \\
\varphi_2)\quad & \beta_1 = b_{1,1}, & \beta_2 = b_{1,2} + b_{2,2}\varTheta \\
\varphi_3)\quad & \gamma_1 = c_{1,1}, & \gamma_2 = c_{1,2} + c_{2,2}\varTheta,
\end{array}
\]
so dass
\[
\omega_1 = \frac{\alpha_2}{\alpha_1}, \quad \omega_2 = \frac{\beta_2}{\beta_1}, \quad \omega_3 = \frac{\gamma_2}{\gamma_1}
\]

wird. Es ist dann $c_{1,1}$ die kleinste positive ganze rationale Zahl,
die durch $\varphi_1\varphi_2$ theilbar ist, und diese muss jedenfalls ein Theiler
von $a_{1,1}b_{1,1}$ sein. Hat man
\[
a_{1,1}\,b_{1,1} = k\,c_{1,1},
\]
wo $k$ eine positive ganze rationale Zahl ist, so ist, weil
$N_a(\varphi_1)\,N_a(\varphi_2) = N_a(\varphi_3)$ sein muss,
\[
k\,a_{2,2}\,b_{2,2} = c_{2,2}.
\]
Die letzte Zahl $c_{1,2}$ ist dann aus der Congruenz
\begin{equation}\label{eq:s168-4}
c_{1,2} + c_{2,2}\varTheta \equiv 0 \pmod{\varphi_1\varphi_2}
\end{equation}
zu bestimmen.

Wenn es nur darauf ankommt, einen Repr"asentanten der
zusammengesetzten Classe zu finden, so verf"ahrt man am einfach-
sten so.

Ist $\varphi_1$ beliebig angenommen, so kann man (nach \S.~152,~6.) $\varphi_2$
in seiner Classe so w"ahlen, dass es relativ prim zu $a_{1,1}$ wird.
Dann ist auch $\varphi_1$ relativ prim zu $\varphi_2$ und $b_{1,1}$ relativ prim zu $a_{1,1}$.

Denn haben $a_{1,1}$ und $b_{1,1}$ einen gr"ossten gemeinschaftlichen
Theiler $d$, so ist dieser auch relativ prim zu $\varphi_2$, und also ist
$b_{1,1} : d$ durch $\varphi_2$ theilbar; dies aber widerspricht, wenn $d > 1$
ist, der Definition von $b_{1,1}$. Ebenso folgt auch umgekehrt, dass,
wenn $a_{1,1}$ und $b_{1,1}$ relativ prim sind, $\varphi_2$ zu $a_{1,1}$ relativ prim ist.
Denn wenn keiner der rationalen Primfactoren von $b_{1,1}$ in $a_{1,1}$
aufgeht, so kann auch keiner der in $\varphi_2$ aufgehenden Prim-
factoren in $a_{1,1}$ aufgehen.

In diesem Falle ist nun:
\begin{equation}\label{eq:s168-5}
a_{1,1}\,b_{1,1} = c_{1,1}, \quad a_{2,2}\,b_{2,2} = c_{2,2}.
\end{equation}
Die Congruenz~(4) wird dann
\begin{equation}\label{eq:s168-6}
c_{1,2} + a_{2,2}b_{2,2}\varTheta \equiv 0 \pmod{\varphi_1\varphi_2},
\end{equation}
und da
\[
a_{1,2} + a_{2,2}\varTheta \equiv 0 \pmod{\varphi_1}, \quad b_{1,2} + b_{2,2}\varTheta \equiv 0 \pmod{\varphi_2}
\]
ist, so zerf"allt~(6) in die beiden Congruenzen:
\begin{alignat*}{2}
c_{1,2} &\equiv b_{2,2}\,a_{1,2} &\pmod{\varphi_1} \\
        &\equiv a_{2,2}\,b_{1,2} &\pmod{\varphi_2}.
\end{alignat*}
Diese Congruenzen sind aber nach \S.~138, 13.\ gleichwerthig mit
den gew"ohnlichen Zahlencongruenzen
\begin{equation}\label{eq:s168-7}
\begin{alignedat}{2}
c_{1,2} &\equiv b_{2,2}\,a_{1,2} &\pmod{a_{1,1}} \\
        &\equiv a_{2,2}\,b_{1,2} &\pmod{b_{1,1}},
\end{alignedat}
\end{equation}
wodurch $c_{1,2}$ nach dem Modul $c_{1,1}$ bestimmt ist.

Es bleibt nur noch "ubrig, zu zeigen, wie man $\varphi_2$ der hier
gestellten Forderung gem"ass bestimmen kann. Da zwei Formen,
die sich nur durch einen Zahlenfactor von einander unter-
scheiden, "aquivalent sind, so ist $\varphi_2$ "aquivalent mit
\[
c_2\,(t_1 + \omega_2t_2),
\]
[\S.~166,~(10)] und nehmen wir diese Form f"ur $\varphi_2$, so ist $b_{1,1} = c_2$.
Die Zahl $\omega_2$ kann man aber unter den mit ihr "aquivalenten
immer so w"ahlen, dass $c_2$ zu $a_{1,1}$ relativ prim wird, und dann ist
auch $a_{1,1}$ zu $\varphi_2$ relativ prim.

Ersetzen wir n"amlich $\omega_2$ durch
\[
\frac{p\,\omega_2 + q}{r\,\omega_2 + s}, \quad ps - qr = \pm 1,
\]
so geht $c_2$ nach Bd.~I, \S.~122,~(13) in
\[
c_2' = -a_2r^2 + b_2rs + c_2s^2
\]
"uber, und nun kann man "uber die ganzen rationalen Zahlen $r,s$
so verf"ugen, dass sie unter einander theilerfremd sind und dass
$c_2'$ durch irgend welche gegebene Primzahlen nicht theilbar ist.

Denn ist $n$ eine solche Primzahl, so nehme man, wenn $n$ in
$a_2$ und $c_2$, also nicht in $b_2$ aufgeht, $r$ und $s$ durch $n$ untheilbar,
wenn $n$ in $a_2$ nicht aufgeht, $s$ durch $n$ theilbar und $r$ durch $n$
untheilbar, und wenn $n$ in $c_2$ nicht aufgeht, $r$ durch $n$ theilbar,
$s$ durch $n$ untheilbar. (Sind $a_2$ und $c_2$ beide nicht durch $n$
theilbar, so hat man zwischen den beiden letzteren Annahmen
die Wahl.) Diese Forderungen sind f"ur verschiedene beliebig
gegebene Primzahlen $n$ mit einander vertr"aglich, und dann l"asst
sich $p, q$ aus der Bedingung $ps - qr = \pm 1$ bestimmen.

Diese Gesetze der Composition der quadratischen Irrational-
zahlen haben wir hier nur unter der Voraussetzung abgeleitet,
dass die Discriminante zugleich die Grundzahl eines quadra-
tischen K"orpers ist, d.~h.\ dass $\varDelta$ Stammdiscriminante ist. Die-
selben Gesetze gelten aber auch, mit geringen Modificationen,
allgemein\footnote{Gauss, Disquisitiones arithm., Art.~234~ff. Dirichlet-Dedekind,
Zahlentheorie. 4.~Auflage, \S.~187.}.

\newpage
\section*{Einundzwanzigster Abschnitt.}
\addcontentsline{toc}{section}{Einundzwanzigster Abschnitt. Kreistheilungsk"orper}

\begin{center}
\Large\bfseries
Kreistheilungsk"orper.
\end{center}
\bigskip

\section*{\S.~169. Zerlegung der Primzahl $q$ in Primfactoren im Kreistheilungsk"orper~$\varOmega_{q^\varkappa}$.}
\addcontentsline{toc}{section}{\S.~169. Zerlegung der Primzahl $q$ in Primfactoren}
\markboth{\S.~169. Zerlegung der Primzahl $q$}{\S.~169. Zerlegung der Primzahl $q$}

Wir machen eine zweite Anwendung der allgemeinen Theorie
der algebraischen Zahlen auf die Kreistheilungsk"orper, und
gewinnen dadurch das Mittel, die Theorie der Abel'schen Zahl-
k"orper "uberhaupt zu einem sch"onen Abschlusse zu bringen.

Die Betrachtungen, die wir zun"achst anzustellen haben,
lassen sich mit kleinen Modificationen auf jeden vollen Kreis-
theilungsk"orper~(\S.~18) anwenden. Der Einfachheit halber be-
schr"anken wir uns hier aber auf den Fall, dass der Grad der
Einheitswurzel eine Primzahlpotenz ist, der f"ur die beabsichtigte
Anwendung ausreicht.

Es sei $q$ eine nat"urliche Primzahl (mit Einschluss von $2$) und
\begin{equation}\label{eq:s169-1}
m = q^\varkappa
\end{equation}
eine Potenz von $q$, deren positiver Exponent $\varkappa$ f"ur $q = 2$ gr"osser
als $1$ vorausgesetzt wird.

Es sei ferner $r$ eine primitive $m^{\mathrm{te}}$ Einheitswurzel und $\varOmega_m$ der
volle Kreistheilungsk"orper f"ur den Exponenten $m$, dessen Grad
\begin{equation}\label{eq:s169-2}
\mu = \varphi(m) = q^{\varkappa-1}(q-1)
\end{equation}
ist (Bd.~I, \S.~134).

Wir bezeichnen durchweg mit $n$ die $\mu$ Zahlen eines vollen
Restsystems f"ur den Modul $m$ mit Ausschluss der durch $q$ theil-
baren Zahlen, und es sind dann die $\mu$ Zahlen $r^n$ die Wurzeln
der irreduciblen Gleichung $\mu^{\mathrm{ten}}$ Grades $f(x) = 0$, worin
\begin{equation}\label{eq:s169-3}
f(x) = \frac{x^m - 1}{x^{\frac{m}{q}} - 1} = x^{q^{\varkappa-1}(q-1)} + x^{q^{\varkappa-1}(q-2)} + \dots + 1
\end{equation}
zu setzen ist. $r$ ist daher eine ganze Zahl in $\varOmega_m$, und ihre
Norm ist $= +1$; folglich ist $r$ eine Einheit.

$\varOmega_m$ ist ein Normalk"orper, der durch die $\mu$ Substitutionen
\[
s_n = (r, r^n)
\]
in sich selber "ubergeht.

Diese $\mu$ Substitutionen $s_n$ bilden eine Abel'sche Gruppe $\mathfrak{A}$,
welche die Galois'sche Gruppe des K"orpers $\varOmega_m$ ist. Sie ist
isomorph mit der gleichfalls durch $\mathfrak{A}$ zu bezeichnenden Gruppe
der nach dem Modul $m$ genommenen Zahlclasses $n$~(\S.~16).

Die Function $f(x)$ l"asst sich in die $\mu$ Factoren $x - r^n$ zer-
legen, und wir setzen daher
\begin{equation}\label{eq:s169-4}
f(x) = \prod_{0, m-1}^{n} (x - r^n).
\end{equation}

Setzen wir darin $x = 1$, und beachten, dass [nach~(3)]
$f(1) = q$ ist, so folgt
\begin{equation}\label{eq:s169-5}
q = \prod_{0, m-1}^{n} (1 - r^n).
\end{equation}
Die Zahlen
\begin{equation}\label{eq:s169-6}
\sigma_n = 1 - r^n, \quad \sigma = 1 - r
\end{equation}
sind aber ganze Zahlen des K"orpers $\varOmega_m$, und die Formel~(5)
zeigt zun"achst, dass $q$ im K"orper $\varOmega_m$ in $\mu$ Factoren $\sigma_n$
zerlegbar ist.

Wir beweisen zun"achst, dass die $\mu$ Zahlen $\sigma_n$ mit einander
associirt sind. Bedeuten $n, n'$ zwei durch $q$ nicht theilbare Zahlen,
so k"onnen wir eine nat"urliche Zahl $a$ so bestimmen, dass
\[
n' \equiv an \pmod{m}
\]
wird. Denn ist aber
\[
\frac{\sigma_{n'}}{\sigma_n} = \frac{1 - r^{an}}{1 - r^n} = 1 + r^n + r^{2n} + \dots + r^{(a-1)n}
\]
eine ganze Zahl, also $\sigma_{n'}$ durch $\sigma_n$ theilbar. Da hierin $n$ mit $n'$
vertauscht werden kann, so ist auch $\sigma_n$ durch $\sigma_{n'}$ theilbar, also
sind beide Zahlen associirt~(\S.~138). Wenn daher $\varepsilon$ eine Einheit
bedeutet, so ist nach~(5):
\begin{equation}\label{eq:s169-7}
q = \varepsilon\,\sigma^\mu, \quad N(\sigma) = q.
\end{equation}

Es ist also die nat"urliche Primzahl $q$ associirt mit der $\mu^{\mathrm{ten}}$
Potenz einer ganzen Zahl $\sigma$ im K"orper $\varOmega_m$. Diese Zahl $\sigma$ ist
aber auch im K"orper $\varOmega_m$ noch Primzahl, und zwar vom ersten
Grade. Denn hat $\sigma$ irgend einen Theiler $\sigma'$, so muss die Norm
von $\sigma'$ ein Theiler der Norm von $\sigma$ sein, also, wenn $\sigma'$ keine

\newpage
\section*{\S.~170. Minimalbasis des K"orpers $\varOmega_m$.}
\addcontentsline{toc}{section}{\S.~170. Minimalbasis des K"orpers $\varOmega_m$}
\markboth{\S.~170. Minimalbasis des K"orpers $\varOmega_m$}{\S.~170. Minimalbasis des K"orpers $\varOmega_m$}

Die S"atze, die im vorigen Paragraphen bewiesen sind, f"uhren
uns zu dem folgenden Theorem:

\begin{theorem}
Die Zahlen
\begin{equation}\label{eq:s170-1}
1,\; r,\; r^2,\; \dots,\; r^{\mu-1}
\end{equation}
bilden eine Basis des Systemes $\mathfrak{o}$ der ganzen
Zahlen in $\varOmega_m$.
\end{theorem}

Um dies zu zeigen, gen"ugt nach \S.~145 der Nachweis, dass
die ganze Zahl in $\varOmega_m$
\begin{equation}\label{eq:s170-2}
\omega = x_0 + x_1r + x_2r^2 + \dots + x_{\mu-1}r^{\mu-1},
\end{equation}
worin $x_0, x_1, \dots, x_{\mu-1}$ ganze rationale Zahlen sind, nur dann
durch eine rationale Primzahl $p$ theilbar sein kann, wenn die
Zahlen $x_0, x_1, \dots, x_{\mu-1}$ alle durch $p$ theilbar sind.

Nehmen wir also an, es sei $\omega$ durch $p$ theilbar; dann sind
auch alle Zahlen $\omega_n$, die aus $\omega$ durch eine der $\mu$ Substitutionen
$(r, r^n)$ entstehen, durch $p$ theilbar, und wir erhalten also aus
(2) ein System von $\mu$ Gleichungen, die in Bezug auf die $x_i$
linear sind:
\begin{equation}\label{eq:s170-3}
\omega_n = x_0 + x_1r^n + x_2r^{2n} + \dots + x_{\mu-1}r^{(\mu-1)n}.
\end{equation}
Die Determinante dieses Systemes, d.~h.\ die aus den $\mu$ Reihen
\[
1,\; r^n,\; r^{2n},\; \dots,\; r^{(\mu-1)n}
\]
gebildete Determinante ist aber nach Bd.~I, \S.~46 gleich der
Quadratwurzel $\sqrt{\varDelta}$, und wenn wir also das System der Glei-
chungen~(3) aufl"osen, so folgt, dass die s"ammtlichen rationalen
Zahlen $\varDelta\,x_i$ durch $p$ theilbar sein m"ussen.

Ist nun $p$ nicht $= q$, so ist $\varDelta$ nicht durch $p$ theilbar, und
s"ammtliche $x_i$ m"ussen durch $p$ theilbar sein.

Ist aber $p = q$, so setzen wir
\[
f(t) = x_0 + x_1t + x_2t^2 + \dots + x_{\mu-1}t^{\mu-1},
\]
so dass $\omega_n = f(r^n)$ wird, und setzen darin $r = 1 - \sigma$ [\S.~169,~(6)].
Dann ist nach der Taylor'schen Entwickelung:
\begin{multline}\label{eq:s170-4}
\omega = f(1-\sigma) = f'(1) - \sigma\,f'(1) + \frac{\sigma^2}{1.2}f''(1) - \dots \\
- \frac{\sigma^{\mu-1}}{\varPi(\mu-1)}f^{(\mu-1)}(1).
\end{multline}
und hierin sind die Coefficienten
\[
f(1),\; f'(1),\; \tfrac{1}{2}f''(1),\; \dots,\; \frac{1}{\varPi(\mu-1)}f^{(\mu-1)}(1)
\]
ganze rationale Zahlen, n"amlich
\begin{equation}\label{eq:s170-5}
\begin{aligned}
\frac{1}{\varPi(\mu-1)}f^{(\mu-1)}(1) &= x_{\mu-1} \\
\frac{1}{\varPi(\mu-2)}f^{(\mu-2)}(1) &= x_{\mu-2} + (\mu-1)x_{\mu-1} \\
\frac{1}{\varPi(\mu-3)}f^{(\mu-3)}(1) &= x_{\mu-3} + (\mu-2)x_{\mu-2} + \frac{(\mu-1)(\mu-2)}{1.2}x_{\mu-1} \\
&\;\vdotswithin{=} \\
f(1) &= x_0 + x_1 + \dots + x_{\mu-1}.
\end{aligned}
\end{equation}
Da nun $\omega$ durch $q$ und folglich durch $\sigma^\mu$ theilbar ist, so
muss $f(1)$ nach~(4) durch $q$ theilbar sein. Da aber $\mu > 1$ und
folglich $q : \sigma$ noch durch $\sigma$ theilbar ist, so folgt, dass auch
\[
f'(1) - \frac{\sigma}{1.2}f''(1) \dots
\]
durch $\sigma$, folglich $f'(1)$ durch $q$ theilbar ist, und wenn man so
weiter schliesst, ergiebt sich, dass die s"ammtlichen rationalen
Zahlen~(5) durch $q$ theilbar sind, und folglich sind auch die
$x_{\mu-1}, x_{\mu-2}, \dots, x_0$ durch $q$ theilbar. Damit ist der Satz~II.\ voll-
st"andig bewiesen.

Anstatt der Reihe~(1) kann man auch eine beliebige Reihe von
$\mu$ auf einander folgenden Potenzen von $r$ als Basis von $\mathfrak{o}$ nehmen:
\begin{equation}\label{eq:s170-6}
r^a,\; r^{a+1},\; r^{a+2},\; \dots,\; r^{a+\mu-1},
\end{equation}
die sich ergiebt, wenn man alle Glieder der Reihe~(1) mit $r^a$
multiplicirt.

Die Richtigkeit dieser Bemerkung ergiebt sich daraus, dass
die $r^{\pm a}$ ganze Zahlen sind, und dass also $\omega$ und $r^a\omega$ stets gleich-
zeitig ganze oder gebrochene Zahlen sind.

Beispielsweise bilden die Zahlen
\begin{equation}\label{eq:s170-7}
r^{-\frac{1}{2}\mu+1},\; r^{-\frac{1}{2}\mu+2},\; \dots,\; r^{-1},\; 1,\; r,\; r^2,\; \dots,\; r^{\frac{1}{2}\mu}
\end{equation}
eine Basis von $\mathfrak{o}$.

Damit ist auch die Grundzahl $\varDelta$ des K"orpers $\varOmega_m$ bestimmt.
Sie ist nach der allgemeinen Definition~(\S.~145) gleich der Dis-
criminante
\[
\varDelta\,(1, r, r^2, \dots, r^{\mu-1}),
\]
d.~h.\ gleich der in \S.~169 bestimmten Discriminante der Kreis-
theilungsgleichung.

\newpage
\section*{\S.~171. Die Primideale im K"orper $\varOmega_m$.}
\addcontentsline{toc}{section}{\S.~171. Die Primideale im K"orper $\varOmega_m$}
\markboth{\S.~171. Die Primideale im K"orper $\varOmega_m$}{\S.~171. Die Primideale im K"orper $\varOmega_m$}

Wir haben im \S.~164 gesehen, dass die nat"urliche Primzahl $q$
die $\mu^{\mathrm{te}}$ Potenz einer Primzahl $\sigma$ im K"orper $\varOmega_m$ ist.

Um alle Primideale, die im K"orper $\varOmega_m$ existiren, zu er-
miteln, sind also noch s"ammtliche von $q$ verschiedene nat"urliche
Primzahlen $p$ in Primfactoren zu zerlegen, und es sind darunter
die nicht associirten auszusuchen.

Die Grundlage f"ur diese Untersuchung bilden die S"atze des
\S.~162, wenn der dort mit $R$ bezeichnete K"orper durch den
K"orper der rationalen Zahlen ersetzt wird.

Jede Zahl $\omega$ des K"orpers $\varOmega_m$ geht durch eine der $\mu$ Sub-
stitutionen
\[
s_n = (r, r^n)
\]
in eine bestimmte andere Zahl $\omega_n$ "uber, die gleichfalls in $\varOmega_m$
enthalten ist. Ist
\begin{equation}\label{eq:s171-1}
\omega = \omega_1 = a_0 + a_1r + a_2r^2 + \dots + a_{\mu-1}r^{\mu-1},
\end{equation}
so ist
\begin{equation}\label{eq:s171-2}
\omega_n = a_0 + a_1r^n + a_2r^{2n} + \dots + a_{\mu-1}r^{(\mu-1)n}.
\end{equation}

Wenn eine dieser Zahlen $\omega_n$ eine ganze Zahl ist, wie wir
jetzt annehmen wollen, so sind es auch alle anderen, und dies
tritt immer dann und nur dann ein, wenn die rationalen Zahlen
$a_0, a_1, \dots, a_{\mu-1}$ ganz sind.

Sind $h, k$ irgend zwei Exponenten, so ist
\[
r^h - r^k = r^k\,(r^{h-k} - 1),
\]
und dies ist nach \S.~169,~(9) mit einer Potenz von $\sigma$ associirt,
also, wenn $\mathfrak{p}$ ein in $p$ aufgehendes Primideal bedeutet, durch $\mathfrak{p}$
nicht theilbar, ausser wenn $h \equiv k \pmod{m}$ ist. Daraus ergiebt
sich, dass zwei Zahlen $\omega_h, \omega_k$ nur dann nach dem Modul $\mathfrak{p}$ con-
gruent sein k"onnen, wenn $h \equiv k \pmod{m}$ ist, und dass also die
Gruppe $X$ des \S.~162, die der Bedingung
\[
\omega \mid \chi \equiv \omega \pmod{\mathfrak{p}}
\]
gen"ugt, nur die identische Substitution enth"alt. Daraus folgt
aber, dass $g = 1$, d.~h.\ dass $p$ nicht durch das Quadrat
eines Primideals theilbar ist.

\newpage
\section*{\S.~172. Die conjugirten Primideale.}
\addcontentsline{toc}{section}{\S.~172. Die conjugirten Primideale}
\markboth{\S.~172. Die conjugirten Primideale}{\S.~172. Die conjugirten Primideale}

Bilden wir aus~(2) die $p^{\mathrm{te}}$ Potenz von $\omega$, und beachten den
Fermat'schen Lehrsatz f"ur rationale Zahlen [$a_i^p \equiv a_i \pmod{p}$],
so ergiebt sich
\[
\omega^p \equiv a_0 + a_1r^p + a_2r^{2p} + \dots + a_{\mu-1}r^{(\mu-1)p} \pmod{p},
\]
oder nach~(2)
\begin{equation}\label{eq:s172-3}
\omega^p \equiv \omega_p \pmod{p}.
\end{equation}

Dadurch ist die Gruppe $\varPsi$ bestimmt, zu der das Ideal $\mathfrak{p}$
geh"ort. Es ist n"amlich nach~(3) auch
\[
\omega^p \equiv \omega_p \pmod{\mathfrak{p}},
\]
daher ist $\psi_0 = (r, r^p)$ und die Gruppe $\varPsi$ besteht nach \S.~162,~(13)
aus den Potenzen dieser Substitution, soweit sie von einander
verschieden sind.

Ist daher $f$ der kleinste positive Exponent, f"ur den
\begin{equation}\label{eq:s172-4}
p^f \equiv 1 \pmod{m}
\end{equation}
ist, d.~h.\ geh"ort $p$ zu dem Exponenten $f$ f"ur den Modul $m$, so ist
\begin{equation}\label{eq:s172-5}
\varPsi = (r, r),\; (r, r^p),\; (r, r^{p^2}),\; \dots,\; (r, r^{p^{f-1}})
\end{equation}
und ist vom $f^{\mathrm{ten}}$ Grade. Demnach ist auch $\mathfrak{p}$ ein Primideal
$f^{\mathrm{ten}}$ Grades und
\begin{equation}\label{eq:s172-6}
N(\mathfrak{p}) = p^f.
\end{equation}

Die Anzahl $e$ der von einander verschiedenen conjugirten
Primfactoren von $p$ ist $\mu : f$, und wir haben also den Satz:

\begin{theorem}
Ist $p$ eine von $q$ verschiedene Primzahl, die f"ur
den Modul $m$ zum Exponenten $f$ geh"ort, ist
ferner $\mu = \varphi(m) = ef$, so zerf"allt $p$ in $\varOmega_m$ in $e$
von einander verschiedene conjugirte Prim-
factoren $f^{\mathrm{ten}}$ Grades.
\end{theorem}

\newpage
Wir m"ussen noch etwas genauer auf die Bildung der con-
jugirten Primfactoren $\mathfrak{p}_n$ einer von $q$ verschiedenen Primzahl $p$
eingehen, die nach dem zuletzt bewiesenen Satze in $\varOmega_m$ in $e$ von
einander verschiedene Primfactoren zerf"allt, wenn
\begin{equation}\label{eq:s172-1}
\varphi(m) = ef
\end{equation}
ist, und $f$ der kleinste positive Exponent, f"ur den
\begin{equation}\label{eq:s172-2}
p^f \equiv 1 \pmod{m}.
\end{equation}

Die Gruppe $\varPsi$ ist isomorph mit einer in $\mathfrak{A}$~(\S.~169) ent-
haltenen Gruppe nach dem Modul $m$ genommener ganzer
rationaler Zahlen, die wir mit $\mathfrak{U}$ bezeichnen, die aus allen, einer
Congruenz
\begin{equation}\label{eq:s172-3b}
a \equiv p^h \pmod{m}
\end{equation}
gen"ugenden Zahlen $a$ besteht, und die wir daher symbolisch durch
\begin{equation}\label{eq:s172-4b}
\mathfrak{U} \equiv p^h \pmod{m}
\end{equation}
darstellen k"onnen, wenn $h$ einen beliebigen ganzzahligen Expo-
nenten bedeutet (vergl.\ \S.~18).

Wir k"onnen also, wenn wir die Zahlen $\xi_1, \xi_2, \dots, \xi_e$ aus $\mathfrak{A}$
passend ausw"ahlen:
\begin{equation}\label{eq:s172-5}
\mathfrak{A} = \mathfrak{U}\xi_1 + \mathfrak{U}\xi_2 + \mathfrak{U}\xi_3 + \dots + \mathfrak{U}\xi_e
\end{equation}
setzen, und jeder dieser Nebengruppen entspricht eines der con-
jugirten Functionale $\mathfrak{p}_{\xi_1}, \mathfrak{p}_{\xi_2}, \dots, \mathfrak{p}_{\xi_e}$, die alle von einander ver-
schieden sind und deren Product mit $p$ associirt ist.

Wir setzen also
\begin{equation}\label{eq:s172-6b}
p = \mathfrak{p}_{\xi_1}\mathfrak{p}_{\xi_2}\dots\mathfrak{p}_{\xi_e}.
\end{equation}

Um aber das Zahlensystem
\begin{equation}\label{eq:s172-7}
\xi_1, \xi_2, \dots, \xi_e,
\end{equation}
auf das es haupts"achlich ankommt, genauer zu charakterisiren,
bemerken wir, dass wir dazu jedes System von $e$ Zahlen der
Gruppe $\mathfrak{A}$ nehmen k"onnen, wenn nur keine zwei unter ihnen,
$\xi, \xi'$, eine Congruenz
\[
\xi'\xi^{-1} \equiv p^h \pmod{m}
\]
erf"ullen.

Um ein solches Zahlensystem zu finden, ist es gut, den Fall
eines ungeraden $m$ von dem anderen zu trennen, in dem $m$ eine
Potenz von $2$ ist.

1. Ist zun"achst $m = q^\varkappa$ ungerade, so nehmen wir eine
primitive Wurzel $g$ von $m$~(\S.~14) und setzen
\begin{equation}\label{eq:s172-8}
p \equiv g^\gamma \pmod{m}.
\end{equation}
Denn $f$ ist [nach~(1),~(2)] die kleinste positive Zahl, die der
Bedingung
\[
\gamma f \equiv 0 \pmod{\mu}
\]
gen"ugt, und folglich ist $e$ der gr"osste gemeinschaftliche Theiler
von $\mu$ und $\gamma$, und irgend eine Zahl $n \equiv g^\lambda$ ist nur dann mit
einer Potenz von $p$ congruent, wenn $\lambda$ durch $e$ theilbar ist. Denn
nur dann kann die Congruenz
\[
\lambda \equiv \gamma x \pmod{\mu}
\]
befriedigt werden. Setzen wir
\[
\gamma = e\,\nu,
\]
so ist $\nu$ relativ prim zu $f$. Daraus folgt, dass die Reihe der Zahlen
\begin{equation}\label{eq:s172-9}
1,\; g,\; g^2,\; \dots,\; g^{e-1}
\end{equation}
f"ur die $\xi$ genommen werden kann. Denn sind $h, h'$ zwei verschiedene
Zahlen der Reihe $0, 1, 2, \dots, e-1$, so ist $h' - h$ niemals durch $e$
theilbar, und folglich
\[
\xi'\xi^{-1} = g^{h'-h}
\]
niemals mit einer Potenz von $p$ congruent.

Wir wollen ferner noch die beiden F"alle unterscheiden, dass
$p - 1$ durch $q$ theilbar ist oder nicht.

a) Ist $p \equiv 1 \pmod{q}$, so ist $\gamma \equiv 0 \pmod{q-1}$, und folg-
lich $e$ theilbar durch $q - 1$. Wir k"onnen daher
\[
e = \varphi(q^{\varkappa_1}) = q^{\varkappa_1-1}(q-1)
\]
setzen, worin $0 < \varkappa_1 \leqq \varkappa$ ist.

Wenn $\varkappa_1 = \varkappa$, also $e = \varphi(m)$, $f = 1$, und daher $p - 1$
durch $m$ theilbar ist, so durchl"auft $\xi$ die ganze Gruppe $\mathfrak{A}$, und
alle Primfunctionale $\mathfrak{p}_\xi$ sind einander verschieden.

Ist $\varkappa_1 < \varkappa$, so ist $\nu$ durch $q$ nicht theilbar, weil sonst $e$ nicht
der gr"osste gemeinschaftliche Theiler von $\mu$ und $\gamma$ w"are, und
\[
p - 1 \equiv g^{\nu\varphi(q^{\varkappa_1})} - 1 \pmod{m}
\]
ist durch $q^{\varkappa_1}$ theilbar, aber durch keine h"ohere Potenz von $q$.
Setzen wir also $q^{\varkappa_1} = m_1$ und
\begin{equation}\label{eq:s172-10}
m = m_1\,m_2,
\end{equation}
so ist $m_1$ der gr"osste gemeinschaftliche Theiler von $p - 1$
und $m$. Die Primzahl $p$ zerf"allt im K"orper $\varOmega_m$ in $\varphi(m_1)$ ver-
schiedene Primfactoren. In ebenso viele Primfactoren zerf"allt
aber $p$ auch im K"orper $\varOmega_{m_1}$, der ein Theiler des K"orpers $\varOmega_m$
ist und aus allen rationalen Functionen der $m_1^{\mathrm{ten}}$ Einheitswurzel
$r_1 = r^{m_2}$ besteht. Daraus folgt, dass die Primfactoren von $p$ in $\varOmega_m$
auch im K"orper $\varOmega_m$ nicht weiter zerlegbar sind, und dass wir
daher die Primfactoren $p$ in $\varOmega_m$ als Functionale des K"orpers
$\varOmega_{m_1}$ darstellen k"onnen.

Es durchl"auft $\xi$ ein volles System durch $q$ nicht theilbarer
Reste nach dem Modul $m_1$, und die $\mathfrak{p}_\xi$ sind Ideale in $\varOmega_{m_1}$.

\newpage
b) Ist umgekehrt $\gamma$ durch $q - 1$ theilbar, so folgt aus~(8),
dass $p - 1$ durch $q$ theilbar ist. Nehmen wir also $p - 1$ nicht
durch $q$ theilbar an, so ist auch $e$ nicht durch $q - 1$ theilbar,
und $\xi$ durchl"auft die Reihe der Zahlen~(9).

Bilden wir die Summe dieser Zahlen $\xi$, so folgt:
\[
(g - 1)\sum\xi \equiv g^e - 1 \pmod{m},
\]
und diese Zahl ist durch $q$ theilbar oder nicht theilbar, je nach-
dem $e$ durch $q - 1$ theilbar oder nicht theilbar ist. Da ausser-
dem $g - 1$ nicht durch $q$ theilbar ist, so ist $\sum\xi$ durch $q$ theilbar
oder nicht theilbar, je nachdem $p - 1$ durch $q$ theilbar oder
nicht theilbar ist.

2. Wir wenden uns zu dem Falle, dass $m = 2^\varkappa$ eine Potenz
von $2$ ist, und hier haben wir wieder zu unterscheiden, ob $p - 1$
durch $4$ oder nur durch $2$ theilbar ist.

a) Ist $p - 1$ durch $4$ theilbar, so ist~(\S.~15):
\[
p \equiv 5^\beta \pmod{m},
\]
und $f$ ist die kleinste positive Zahl, die der Bedingung
\[
\beta f \equiv 0 \;\left[\bmod \tfrac{1}{2}\varphi(m)\right]
\]
gen"ugt. Ist $f = 1$, also $\beta$ durch $\tfrac{1}{2}\varphi(m) = 2^{\varkappa-2}$ theilbar, so
ist $e = \varphi(m)$, und $\xi$ durchl"auft die ganze Gruppe $\mathfrak{A}$, d.~h.\ alle
ungeraden Zahlen zwischen $0$ und $m$. Dies ist der Fall, wo
$p - 1$ durch $m$ theilbar ist.

Ist $2^{\varkappa_1-2}$ die h"ochste Potenz von $2$, die in $\beta$ aufgeht, und
$2 \leqq \varkappa_1 < \varkappa$, so ist, wenn wir $m_1 = 2^{\varkappa_1}$ und $m = m_1m_2$ setzen,
\[
f = 2^{\varkappa-\varkappa_1}, \quad e = 2^{\varkappa_1-1} = \varphi(m_1),
\]
und es ist, wenn $s$ eine ungerade Zahl ist,
\[
p - 1 \equiv 5^{s\cdot\frac{1}{2}\varphi(m_1)} - 1 \dots \pmod{m}.
\]
Nach \S.~15 ist $5^{\frac{1}{2}\varphi(m_1)} - 1$ durch $m_1$, aber durch keine
h"ohere Potenz von $2$ theilbar. Da man ausserdem nach den-
selben S"atzen
\[
5^{s\cdot\frac{1}{2}\varphi(m_1)} \equiv 5^{\frac{1}{2}\varphi(m_1)} \pmod{2m_1}
\]
hat, so folgt, dass $m_1$ der gr"osste gemeinschaftliche Theiler von
$p - 1$ und $m$ ist.

Wenn wir also, wie oben, den K"orper $\varOmega_{m_1}$ zuziehen, so er-
giebt sich, dass, wenn $m_1 \geqq 4$ der gr"osste gemeinschaftliche
Theiler von $p - 1$ und $m$ ist, $\xi$ ein volles System ungerader
Reste von $m_1$ durchl"auft, und dass $\mathfrak{p}_\xi$ zugleich die Primfactoren
von $p$ im K"orper $\varOmega_{m_1}$ sind.

b) Ist $p - 1$ nicht durch $4$ theilbar, also
\begin{equation}\label{eq:s172-11}
p \equiv -5^\beta \pmod{m},
\end{equation}
so ist $f$ die kleinste positive Zahl, die den Congruenzen
\[
f \equiv 0 \pmod{2}, \quad 2\beta f \equiv 0 \;\left[\bmod \varphi(m)\right]
\]
gen"ugt, und folglich ist $e$ der gr"osste gemeinschaftliche Theiler
von $\varphi(m)$ und $2\beta$, oder, wenn $2\beta$ durch $\varphi(m)$ theilbar ist,
gleich $\tfrac{1}{2}\varphi(m)$. (Eine Ausnahme bildet hier der Fall $m = 4$, in
dem $e = 1$ ist. Diesen Fall, der ja schon in fr"uheren Abschnitten
vollst"andig erledigt ist, lassen wir daher jetzt bei Seite.)

Es ergiebt sich dann aus~(11), dass eine Zahl von der Form $5^\lambda$
nur dann einer Potenz von $p$ congruent sein kann, wenn $\lambda$ ein
Vielfaches von $2\beta$, also ein Vielfaches von $e$ ist.

Hier kann nun in jeder der Nebengruppen $\mathfrak{U}\xi$ eine Zahl
gefunden werden, die $\equiv 1 \pmod{4}$ ist, da $p \equiv -1 \pmod{4}$
eine Zahl in $\mathfrak{A}$ ist, und wir erhalten demnach ein vollst"andiges
System der $\xi$ in der Form:
\[
1,\; 5,\; 5^2,\; \dots,\; 5^{e-1},
\]
und darin ist $e$ h"ochstens $= \tfrac{1}{2}\varphi(m) = 2^{\varkappa-2}$.

Wir fassen das Bewiesene im folgenden Theorem zusammen:

\begin{theorem}
Zerlegt man die von $q$ verschiedene Primzahl $p$
im K"orper $\varOmega_m$ in ihre Primfactoren:
\[
p = \mathfrak{p}_{\xi_1}\mathfrak{p}_{\xi_2}\dots\mathfrak{p}_{\xi_e},
\]
so kann man, wenn $p - 1$ durch $q$, oder im Falle
eines geraden $m$ durch $4$ theilbar ist, die $\xi$ ein
volles System durch $q$ nicht theilbarer Reste
nach dem Modul $m_1$ durchlaufen lassen, und
die $\mathfrak{p}_\xi$ als Primfunctionale im K"orper $\varOmega_{m_1}$ an-
nehmen, wenn $m_1$ der gr"osste gemeinschaftliche
Theiler von $p - 1$ und $m$ ist.

Ist aber $p - 1$ nicht durch $q$ oder f"ur ein
gerades $m$ nicht durch $4$ theilbar, und bedeutet $g$
im ersten Falle eine Primitivwurzel von $m$, im
zweiten Falle die Zahl $5$, so durchl"auft $\xi$ die
Reihe der Zahlen
\[
1,\; g,\; g^2,\; \dots,\; g^{e-1},
\]
und $e$ ist bei einem ungeraden $m$ nicht durch
$q - 1$ theilbar, und bei geradem $m$ mindestens
$= 2$ und h"ochstens $= \tfrac{1}{2}\varphi(m)$. Ausgeschlossen ist
hierbei der Fall $m = 4$.
\end{theorem}

\newpage
\section*{\S.~173. Darstellung der Primfactoren von $p$.}
\addcontentsline{toc}{section}{\S.~173. Darstellung der Primfactoren von $p$}
\markboth{\S.~173. Darstellung der Primfactoren von $p$}{\S.~173. Darstellung der Primfactoren von $p$}

Zur Darstellung der Primfunctionale des K"orpers $\varOmega_m$ k"onnen
wir ein Verfahren anwenden, was wir im \S.~157 kennen gelernt
haben, welches sich hier in Folge des Umstandes, dass
\[
1,\; r,\; r^2,\; \dots,\; r^{\mu-1}
\]
eine Basis von $\mathfrak{o}$ ist, wesentlich vereinfacht.

Die nat"urliche Primzahl $q$ ist, wie wir schon gesehen haben,
die $\mu^{\mathrm{te}}$ Potenz einer im K"orper $\varOmega_m$ existirenden Primzahl $\sigma$; mit
dieser brauchen wir uns nicht weiter zu besch"aftigen, und be-
trachten daher hier nur die von $q$ verschiedenen Primzahlen~$p$.
Es sei $\mathfrak{p}$ ein Primfactor von $p$, und $\mathfrak{U}$ habe dieselbe Bedeutung
wie oben. Wenn $\mathfrak{p}$ zum Exponenten $f$ geh"ort und
\[
ef = \varphi(m) = \mu
\]
ist, so besteht die Gruppe $\mathfrak{U}$ aus den Zahlen
\[
1,\; p,\; p^2,\; \dots,\; p^{f-1}.
\]

Wenn nun
\[
f(t) = N(t-r) = \prod_{0, m-1}^{n}(t - r^n)
\]
das Polynom von $t$ vom Grade $\mu$ bedeutet, dessen Wurzeln die
$\mu$ Gr"ossen $r^n$ sind, so k"onnen wir dies so in Factoren zerlegen,
dass jeder Factor einer der Nebengruppen von $\mathfrak{U}$ entspricht.
Setzen wir n"amlich
\begin{equation}\label{eq:s173-1}
F_1(t) = \prod_{0, f-1}^{h}(t - r^{p^h}),
\end{equation}
und allgemein f"ur jedes durch $q$ nicht theilbare $n$
\begin{equation}\label{eq:s173-2}
F_n(t) = \prod_{0, f-1}^{h}(t - r^{np^h}),
\end{equation}
so ist $F_n(t)$ mit $F_{n'}(t)$ identisch, wenn $n$ und $n'$ in dieselbe
Nebengruppe $\mathfrak{U}\xi_i$ geh"oren. Wenn wir also mit $\xi_1, \xi_2, \dots, \xi_e$ das
im \S.~172,~(5) angewandte Zahlensystem verstehen, so ist
\begin{equation}\label{eq:s173-3}
f(t) = F_{\xi_1}(t)\,F_{\xi_2}(t)\dots F_{\xi_e}(t).
\end{equation}

Nun ist aber nach dem binomischen Lehrsatze:
\[
(t - r^{p^h})^p \equiv (t^p - r^{p^{h+1}}) \pmod{p},
\]
\end{document}
